% ============================================================
%  Forces and Motion — Lesson Booklet (one section per lesson: A Reading · B Theory · C Practice · D Above level)
%  Years 7–8  |  Victorian Curriculum F–10 Version 2.0
%  Compile with: pdflatex booklet.tex  (twice for page refs)
% ============================================================
\documentclass[11pt,a4paper]{article}

% ── Fonts (compile with LuaLaTeX) ────────────────────────────────────────────
\usepackage{fontspec}
% Bundled fonts (in ./fonts/) so the document compiles anywhere with XeLaTeX/LuaLaTeX
\setmainfont{Lato}[                   % body — humanist sans (stands in for deck's Inter)
  Path=./fonts/, Extension=.ttf, Ligatures=TeX,
  UprightFont=Lato-Regular, ItalicFont=Lato-Italic,
  BoldFont=Lato-Bold, BoldItalicFont=Lato-BoldItalic ]
\newfontfamily\headingfont{Poppins}[  % display/headings — matches the slide decks
  Path=./fonts/, Extension=.ttf, Ligatures=TeX,
  UprightFont=Poppins-Regular, ItalicFont=Poppins-Italic,
  BoldFont=Poppins-Bold, BoldItalicFont=Poppins-BoldItalic ]
\setlength{\headheight}{22pt}

\usepackage[protrusion=true,expansion=false]{microtype}
\usepackage[a4paper, top=2cm, bottom=2cm, left=2cm, right=2cm]{geometry}
\usepackage{xcolor}
\usepackage{tikz}
\usetikzlibrary{arrows.meta, shapes.geometric, calc, positioning, decorations.pathreplacing}
\usepackage{tcolorbox}
\tcbuselibrary{skins, breakable, raster}
\usepackage{fancyhdr}
\usepackage{titlesec}
\usepackage{enumitem}
\usepackage{booktabs}
\usepackage{tabularx}
\usepackage{array}
\usepackage{colortbl}
\usepackage{multicol}
\usepackage{amsmath}
\usepackage{needspace}
% TikZ diagram macros are embedded below (just before \begin{document}) for easy editing.
\usepackage{longtable}
\usepackage{lastpage}
\usepackage{calc}
\usepackage{mdframed}
\usepackage{setspace}
\usepackage{ragged2e}

% ── Colour Palette ───────────────────────────────────────────────────────────
% Remapped to the slide-deck palette (names kept so body + diagrams cascade)
\definecolor{NavyBlue}   {HTML}{16305C}
\definecolor{MidBlue}    {HTML}{2E6FE0}
\definecolor{LightBlue}  {HTML}{E8F1FF}
\definecolor{TealGreen}  {HTML}{0F6E63}
\definecolor{LightGreen} {HTML}{E0F7F3}
\definecolor{GoldDark}   {HTML}{C77F12}
\definecolor{GoldLight}  {HTML}{FFF3DC}
\definecolor{PurpleDark} {HTML}{5B3FC0}
\definecolor{PurpleLight}{HTML}{EFE9FF}
\definecolor{OrangeDark} {HTML}{C9483D}
\definecolor{OrangeLight}{HTML}{FFE9E6}
\definecolor{MidGrey}    {HTML}{AEBED4}
\definecolor{LightGrey}  {HTML}{F2F5FA}
\definecolor{DarkGrey}   {HTML}{3A4759}
\definecolor{TableHead}  {HTML}{16305C}
\definecolor{RowAlt}     {HTML}{EEF4FF}
% Bright accent tints (deck) for bars, rules and chips
\definecolor{Blue}   {HTML}{2E6FE0}
\definecolor{Teal}   {HTML}{18A999}
\definecolor{Amber}  {HTML}{F5A623}
\definecolor{Coral}  {HTML}{FF6B5E}
\definecolor{Violet} {HTML}{7C5CE0}
\definecolor{Sky}    {HTML}{4DA3FF}
\definecolor{Ink}    {HTML}{16223A}
\definecolor{Line}   {HTML}{DDE5F0}
\definecolor{Paper}  {HTML}{FBFCFE}

% ── Page Style ───────────────────────────────────────────────────────────────
\pagestyle{fancy}
\fancyhf{}
\renewcommand{\headrulewidth}{0.4pt}
\renewcommand{\headrule}{\hbox to\headwidth{\color{NavyBlue}\leaders\hrule height \headrulewidth\hfill}}
\fancyhead[L]{\headingfont\small\color{NavyBlue}Forces and Motion}
\fancyhead[R]{\headingfont\small\color{DarkGrey}Years 7--8 \textbar{} Victorian Curriculum F--10 v2.0}
\fancyfoot[C]{\headingfont\small\color{DarkGrey}\thepage\ \textbar{} \pageref{LastPage}}
\fancypagestyle{plain}{%
  \fancyhf{}
  \fancyfoot[C]{\headingfont\small\color{DarkGrey}\thepage\ \textbar{} \pageref{LastPage}}
  \renewcommand{\headrulewidth}{0pt}
}

% ── Section Formatting ───────────────────────────────────────────────────────
\titleformat{\section}
  {\headingfont\color{NavyBlue}\Large\bfseries}
  {}{0em}{}
  [{\color{Teal}\rule[0.6ex]{2.2em}{3pt}}{\color{Line}\rule[0.6ex]{\dimexpr\linewidth-2.2em\relax}{1pt}}]
\titlespacing{\section}{0pt}{18pt}{8pt}

\titleformat{\subsection}
  {\headingfont\color{MidBlue}\large\bfseries}
  {\thesubsection}{0.6em}{}
\titlespacing{\subsection}{0pt}{12pt}{4pt}

\titleformat{\subsubsection}
  {\headingfont\color{NavyBlue}\normalsize\bfseries}
  {}{0em}{}
\titlespacing{\subsubsection}{0pt}{8pt}{2pt}

% Deck-style eyebrow / kicker (optional use)
\newcommand{\kicker}[2][Teal]{%
  {\headingfont\bfseries\footnotesize\color{#1}%
   \textcolor{#1}{\rule[0.18ex]{1.6em}{2.2pt}}\hspace{0.5em}\MakeUppercase{#2}\par}%
  \vspace{2pt}}

% ── Custom tcolorbox Environments ────────────────────────────────────────────

% Section divider banner
\newtcolorbox{sectionbanner}[2]{
  enhanced, breakable=false,
  colback=NavyBlue, colframe=NavyBlue,
  fontupper=\color{white}\bfseries\large,
  left=4mm, right=4mm, top=3mm, bottom=3mm,
  title={\color{white}\Huge\bfseries #1},
  attach boxed title to top left={yshift=-2mm, xshift=4mm},
  boxed title style={colback=NavyBlue, colframe=NavyBlue},
  after upper={\par\smallskip\normalfont\normalsize\color{white!80!NavyBlue}\textit{#2}}
}

% Key info box (blue)
\newtcolorbox{infobox}[1]{
  breakable, arc=2.5mm, boxrule=0.5pt,
  colback=LightBlue, colframe=Blue!40,
  fonttitle=\headingfont\bfseries, coltitle=MidBlue,
  title={#1},
  left=4mm, right=3.5mm, top=2.5mm, bottom=2.5mm
}

% Scaffold box (yellow)
\newtcolorbox{scaffold}[1]{
  breakable, arc=2.5mm, boxrule=0.5pt,
  colback=GoldLight, colframe=Amber!40,
  fonttitle=\headingfont\bfseries, coltitle=GoldDark,
  title={#1},
  left=4mm, right=3.5mm, top=2.5mm, bottom=2.5mm
}

% Worked example box
\newtcolorbox{workedexbox}[1]{
  breakable, arc=2.5mm, boxrule=0.5pt,
  colback=white, colframe=Amber!50,
  fonttitle=\headingfont\bfseries, coltitle=GoldDark,
  title={#1},
  left=4mm, right=3.5mm, top=2.5mm, bottom=2.5mm
}
% Worked examples: \begin{workedex}{title} -> "Worked Example N.M --- title",
% auto-numbered per section and labelled we:N.M. Never type the number.
\newcounter{wenum}[section]
\renewcommand{\thewenum}{\thesection.\arabic{wenum}}
\newenvironment{workedex}[1]
  {\refstepcounter{wenum}\label{we:\thewenum}\begin{workedexbox}{Worked Example \thewenum\ --- #1}}
  {\end{workedexbox}}

% Year 7 box (blue)
\newtcolorbox{yr7box}{
  breakable, arc=2.5mm, boxrule=0.5pt,
  colback=LightBlue, colframe=Blue!40,
  left=4mm, right=3.5mm, top=2.5mm, bottom=2.5mm
}

% Year 8 extension box (violet)
\newtcolorbox{yr8box}{
  breakable, arc=2.5mm, boxrule=0.5pt,
  colback=PurpleLight, colframe=Violet!40,
  left=4mm, right=3.5mm, top=2.5mm, bottom=2.5mm
}

% Answer lines box
\newcommand{\answerbox}[1]{%
  \begin{tcolorbox}[
    enhanced, breakable, arc=2mm,
    colback=Paper, colframe=Line,
    left=2mm, right=2mm, top=1mm, bottom=1mm,
    boxrule=1pt,
    height=#1
  ]
  \end{tcolorbox}%
}

% Experiment header
\newtcolorbox{expheader}[3]{
  breakable=false, arc=2.5mm,
  colback=RowAlt, colframe=Amber!50, boxrule=0.5pt,
  fonttitle=\headingfont\bfseries\large, coltitle=NavyBlue,
  title={Experiment #1: #2},
  left=4mm, right=3.5mm, top=2.5mm, bottom=2.5mm,
  after upper={\par\smallskip
    {\headingfont\bfseries\color{GoldDark}Aim:}\ \textit{#3}},
}

% Reading comprehension passage box (green)
\newtcolorbox{readingbox}[1]{
  breakable, arc=2.5mm, boxrule=0.5pt,
  colback=white, colframe=Teal!50,
  fonttitle=\headingfont\bfseries, coltitle=TealGreen,
  title={#1},
  left=4mm, right=3.5mm, top=2.5mm, bottom=2.5mm,
  fontupper=\normalsize
}

% Skill-builder box (for beginner scaffolding)
\newtcolorbox{skillbox}[1]{
  breakable, arc=2.5mm, boxrule=0.5pt,
  colback=RowAlt, colframe=Blue!40,
  fonttitle=\headingfont\bfseries, coltitle=MidBlue,
  title={#1},
  left=4mm, right=3.5mm, top=2.5mm, bottom=2.5mm
}

% ── Custom Commands ───────────────────────────────────────────────────────────

% Question: \question{marks}{text}  (marks can be empty)
% Numbered automatically as Q<section>.<n>, reset every section.
% Do NOT write numbers by hand.
\newcounter{qnum}[section]
\renewcommand{\theqnum}{\thesection.\arabic{qnum}}
\newcommand{\question}[2]{%
  \refstepcounter{qnum}%
  \needspace{3\baselineskip}%
  \medskip
  \noindent\label{q:\theqnum}\textbf{Q\theqnum}\ifx&#1&\else\ \textcolor{DarkGrey}{(#1 mark\ifnum#1=1\else s\fi)}\fi\quad #2%
  \medskip
}

% Inline marks note
\newcommand{\markscount}[1]{\hfill\textcolor{DarkGrey}{\textit{(#1 mark\ifnum#1=1\else s\fi)}}}

% ── Practice set counter and question macros ─────────────────────────────────
\newcounter{pqnum}
\newcommand{\practiceset}[1]{%
  \needspace{6\baselineskip}
  \setcounter{pqnum}{0}%
  \vspace{6pt}
  {\headingfont\bfseries\large\color{GoldDark} Practice Set --- #1}\\[-2pt]
  {\color{Amber}\rule{\linewidth}{1pt}}\\[2pt]
}
% \pq{question text}{working box height}
\newcommand{\pq}[2]{%
  \stepcounter{pqnum}%
  \needspace{4\baselineskip}
  \vspace{5pt}
  \noindent\textbf{\thepqnum.}\quad #1\par
  \vspace{2pt}
  \answerbox{#2}
}
% Two short questions side by side, each with its own box
\newcommand{\pqpair}[4]{%
  \needspace{4\baselineskip}
  \vspace{5pt}
  \noindent
  \begin{minipage}[t]{0.48\linewidth}
    \stepcounter{pqnum}\textbf{\thepqnum.}\quad #1\par\vspace{2pt}
    \begin{tcolorbox}[enhanced,arc=2mm,colback=Paper,colframe=Line,left=2mm,right=2mm,
      top=1mm,bottom=1mm,boxrule=1pt,height=#2]\end{tcolorbox}
  \end{minipage}\hfill
  \begin{minipage}[t]{0.48\linewidth}
    \stepcounter{pqnum}\textbf{\thepqnum.}\quad #3\par\vspace{2pt}
    \begin{tcolorbox}[enhanced,arc=2mm,colback=Paper,colframe=Line,left=2mm,right=2mm,
      top=1mm,bottom=1mm,boxrule=1pt,height=#4]\end{tcolorbox}
  \end{minipage}
}

% ── Comprehension question macro (compact, lined) ────────────────────────────
\newcounter{rcnum}
\newcommand{\rcreset}{\setcounter{rcnum}{0}}
\newcommand{\rcq}[2]{%
  \stepcounter{rcnum}%
  \needspace{3\baselineskip}
  \vspace{4pt}
  \noindent\textbf{C\thercnum.}\quad #1\par
  \vspace{2pt}
  \answerbox{#2}
}

% Section divider (big numbered banner): \sectiondiv{title}{summary}
%
% This IS the section: it steps the real section counter, so subsections,
% tasks, questions and figures inside it are numbered N.M automatically.
% Put \label{sec:N} straight after it. Never hand-type a section number.
\newcommand{\sectiondiv}[2]{%
  \needspace{5\baselineskip}%
  \refstepcounter{section}%
  \begin{tcolorbox}[
    enhanced, breakable=false, arc=3.5mm,
    colback=NavyBlue, colframe=NavyBlue, boxrule=0pt,
    sidebyside, sidebyside align=center,
    lefthand width=2.0cm,
    left=2mm, right=4mm, top=4mm, bottom=4mm
  ]
    \begin{center}
      {\headingfont\bfseries\color{white!35!NavyBlue}\fontsize{50}{50}\selectfont \thesection}
    \end{center}
    \tcblower
    {\headingfont\bfseries\color{white}\LARGE #1}\\[3pt]
    {\itshape\color{white!75!NavyBlue}\small #2}
  \end{tcolorbox}%
}

% Tasks: \task, \task[a] (part a of a NEW task), \taskpart{b} (next part
% of the same task), \exttask (Year 8 "Extension" — same counter).
% Numbered Task <section>.<n>, reset every section. Do NOT type numbers.
% Every task, question and figure also gets an automatic label
% (task:1.7b, q:2.4, fig:1.2, fig:E2.1) so page refs come from the .aux.
\newcounter{tasknum}[section]
\renewcommand{\thetasknum}{\thesection.\arabic{tasknum}}
\NewDocumentCommand{\task}{o}{\refstepcounter{tasknum}\label{task:\thetasknum\IfValueT{#1}{#1}}\textbf{Task \thetasknum\IfValueT{#1}{#1}}}
\newcommand{\taskpart}[1]{\label{task:\thetasknum#1}\textbf{Task \thetasknum#1}}
\newcommand{\exttask}{\refstepcounter{tasknum}\label{task:\thetasknum}\textbf{Extension \thetasknum}}

% Figure captions: \figcap[fig:label]{caption} -> "Figure N.M --- caption"
% \expfigcap{caption} -> "Figure EN.M" (experiment pages). Both auto-numbered.
\newcounter{fignum}[section]
\renewcommand{\thefignum}{\thesection.\arabic{fignum}}
\newcounter{expfignum}[section]
\renewcommand{\theexpfignum}{E\thesection.\arabic{expfignum}}
\newcommand{\figcap}[2][]{%
  \begin{center}
    \refstepcounter{fignum}\label{fig:\thefignum}\if\relax\detokenize{#1}\relax\else\label{#1}\fi
    \small\itshape\color{DarkGrey} Figure \thefignum\ --- #2
  \end{center}%
}
\newcommand{\expfigcap}[1]{%
  \begin{center}
    \refstepcounter{expfignum}\label{fig:\theexpfignum}%
    \small\itshape\color{DarkGrey} Figure \theexpfignum\ --- #1
  \end{center}%
}

% Diagram include
\newcommand{\diagramfig}[3]{%  {filename}{width fraction}{caption}
  \begin{center}
    \includegraphics[width=#2\linewidth]{#1}\\[4pt]
    \small\itshape\color{DarkGrey} #3
  \end{center}%
}

% Table column types
\newcolumntype{L}[1]{>{\RaggedRight\arraybackslash}p{#1}}
\newcolumntype{C}[1]{>{\Centering\arraybackslash}p{#1}}

% Table header colour
\newcommand{\thead}[1]{\cellcolor{TableHead}\color{white}\bfseries #1}
\newcommand{\talt}{\rowcolor{RowAlt}}
\newcommand{\tmid}{\rowcolor{MidGrey}}

% Inline year tags
\newcommand{\tagyr}[1]{%
  \tikz[baseline=-0.5ex]\node[
    rounded corners=2pt,
    fill=#1!20, draw=#1, text=#1,
    inner xsep=4pt, inner ysep=1pt,
    font=\scriptsize\bfseries
  ]{Yr #1 only};}

% Draw answer lines
\newcommand{\anslines}[1]{%
  \vspace{2pt}%
  \foreach \n in {1,...,#1}{%
    \par\noindent\rule{\linewidth}{0.3pt}\vspace{3pt}%
  }%
}

% ── Document ─────────────────────────────────────────────────────────────────
% ══════════════════════════════════════════════════════════════════════════════
%  EMBEDDED DIAGRAM DEFINITIONS  (formerly tikz_diagrams.tex)
%  Edit any figure here; each \diag... macro is called from the body by name.
% ══════════════════════════════════════════════════════════════════════════════
% ============================================================
%  TikZ Diagram Definitions — Forces and Motion Unit Booklet
%  Self-contained macros, no external image files.
%  Requires: tikz with arrows.meta, calc, positioning, patterns
% ============================================================

\usetikzlibrary{arrows.meta,calc,positioning,patterns,backgrounds}

\tikzset{
  force arrow/.style={-{Stealth[length=7pt,width=5pt]}, line width=2pt},
  dim arrow/.style={<->, >=Stealth, thin, color=DarkGrey},
  thin arrow/.style={-{Stealth[length=6pt,width=4pt]}, line width=1.4pt},
  ground hatch/.style={pattern=north east lines, pattern color=DarkGrey!50},
}

% ── D1: Contact vs Non-contact Forces ─────────────────────────────────────────
\newcommand{\diagForceTypes}{%
\begin{center}
\begin{tikzpicture}[font=\small]
  % Contact forces box
  \node[draw=MidBlue, fill=LightBlue, rounded corners=5pt,
        minimum width=5.6cm, minimum height=3.8cm] (cbox) at (0,0) {};
  \node[font=\small\bfseries, text=NavyBlue] at (0, 1.65) {Contact Forces};
  \node[anchor=west, text=NavyBlue, font=\small] at (-2.5, 1.1)  {\textbullet\ Friction};
  \node[anchor=west, text=NavyBlue, font=\small] at (-2.5, 0.65) {\textbullet\ Normal force};
  \node[anchor=west, text=NavyBlue, font=\small] at (-2.5, 0.20) {\textbullet\ Tension};
  \node[anchor=west, text=NavyBlue, font=\small] at (-2.5,-0.25) {\textbullet\ Applied force};
  \node[anchor=west, text=NavyBlue, font=\small] at (-2.5,-0.70) {\textbullet\ Air resistance};
  \node[anchor=west, text=NavyBlue, font=\small] at (-2.5,-1.15) {\textbullet\ Spring force};
  % Non-contact forces box
  \node[draw=red!55!black, fill=red!8, rounded corners=5pt,
        minimum width=5.6cm, minimum height=3.8cm] (ncbox) at (6.4,0) {};
  \node[font=\small\bfseries, text=red!55!black] at (6.4, 1.65) {Non-contact Forces};
  \node[anchor=west, text=red!55!black, font=\small] at (3.9, 1.1)  {\textbullet\ Gravity};
  \node[anchor=west, text=red!55!black, font=\small] at (3.9, 0.65) {\textbullet\ Magnetic force};
  \node[anchor=west, text=red!55!black, font=\small] at (3.9, 0.20) {\textbullet\ Electrostatic force};
  \node[anchor=west, text=red!55!black, font=\small] at (3.9,-0.25) {\textbullet\ Nuclear force};
\end{tikzpicture}
\end{center}}

% ── D2: Free Body Diagrams ────────────────────────────────────────────────────
\newcommand{\diagFBD}{%
\begin{center}
\begin{tikzpicture}[font=\small]
  % ── LEFT: Balanced ──
  \node[font=\small\bfseries, text=NavyBlue] at (0,3.6) {Balanced Forces};
  \fill[ground hatch] (-1.5,-1.0) rectangle (1.5,-0.7);
  \draw[thick, DarkGrey] (-1.5,-0.7) -- (1.5,-0.7);
  \node[draw=MidBlue, fill=LightBlue, rounded corners=3pt,
        minimum width=1.8cm, minimum height=1.0cm,
        font=\small\bfseries, text=NavyBlue] (book) at (0,0) {Book};
  % Normal up
  \draw[force arrow, TealGreen] (book.north) -- ++(0,1.8)
        node[right, font=\small\bfseries, text=TealGreen] {$N=10$\,N};
  % Weight down
  \draw[force arrow, red!65!black] (book.south) -- ++(0,-1.8)
        node[right, font=\small\bfseries, text=red!65!black] {$W=10$\,N};
  \node[font=\footnotesize\itshape, text=DarkGrey] at (0,-3.2) {Net force $= 0$\,N};
  \node[font=\footnotesize\itshape, text=DarkGrey] at (0,-3.65) {No acceleration};

  % ── RIGHT: Unbalanced ──
  \node[font=\small\bfseries, text=NavyBlue] at (7.5,3.6) {Unbalanced Forces};
  \fill[ground hatch] (6.0,-1.0) rectangle (9.0,-0.7);
  \draw[thick, DarkGrey] (6.0,-0.7) -- (9.0,-0.7);
  \node[draw=MidBlue, fill=LightBlue, rounded corners=3pt,
        minimum width=1.8cm, minimum height=1.0cm,
        font=\small\bfseries, text=NavyBlue] (box) at (7.5,0) {Box};
  % Normal up
  \draw[force arrow, TealGreen] (box.north) -- ++(0,1.8)
        node[right, font=\small\bfseries, text=TealGreen] {$N$};
  % Weight down
  \draw[force arrow, red!65!black] (box.south) -- ++(0,-1.8)
        node[right, font=\small\bfseries, text=red!65!black] {$W$};
  % Push right (20N, longer)
  \draw[force arrow, PurpleDark] (box.east) -- ++(2.6,0)
        node[above, font=\small\bfseries, text=PurpleDark] {Push 20\,N};
  % Friction left (5N, shorter)
  \draw[force arrow, OrangeDark] (box.west) -- ++(-1.6,0)
        node[above, font=\small\bfseries, text=OrangeDark] {Friction 5\,N};
  \node[font=\footnotesize\itshape, text=DarkGrey] at (7.5,-3.2) {Net force $= 15$\,N $\rightarrow$};
  \node[font=\footnotesize\itshape, text=DarkGrey] at (7.5,-3.65) {Accelerates rightward};
\end{tikzpicture}
\end{center}}

% ── D3: Newton's First Law ─────────────────────────────────────────────────────
\newcommand{\diagNewtonOne}{%
\begin{center}
\begin{tikzpicture}[font=\small]
  % Law statement
  \node[draw=PurpleDark, fill=PurpleLight!50, rounded corners=5pt,
        minimum width=13cm, minimum height=1.1cm, align=center,
        font=\small\bfseries, text=PurpleDark] at (5.5,2.8)
        {\textit{``An object at rest stays at rest, and an object in motion stays in}\\[-1pt]
         \textit{motion, unless acted on by an unbalanced force.''\ (Law of Inertia)}};

  % Ground lines
  \draw[thick, DarkGrey] (-0.4,0) -- (2.2,0);
  \draw[thick, DarkGrey] (3.6,0) -- (7.0,0);
  \draw[thick, DarkGrey] (8.5,0) -- (11.5,0);

  % Ball 1: at rest
  \shade[ball color=MidBlue!80] (0.9,0.42) circle (0.42);
  \node[font=\footnotesize\bfseries, text=NavyBlue] at (0.9,1.9) {At rest};
  \draw[force arrow, TealGreen]   (0.9,0.84) -- (0.9,1.5);
  \draw[force arrow, red!65!black](0.9,0.0)  -- (0.9,-0.65);
  \node[font=\scriptsize, text=TealGreen, right]    at (0.9,1.2)  {$N$};
  \node[font=\scriptsize, text=red!65!black, right] at (0.9,-0.3) {$W$};
  \node[font=\scriptsize\itshape, text=DarkGrey] at (0.9,-1.1) {stays at rest};

  % Arrow 1→2
  \draw[-{Stealth}, thick, DarkGrey!60] (2.4,-2) -- (3.4,-2)
        node[midway, above, font=\scriptsize\itshape, text=DarkGrey] {no force added};

  % Ball 2: rolling
  \shade[ball color=TealGreen!70] (5.3,0.42) circle (0.42);
  \node[font=\footnotesize\bfseries, text=TealGreen] at (5.3,1.9) {Rolling (no friction)};
  \draw[force arrow, MidBlue] (5.72,0.42) -- (7.0,0.42)
        node[above, font=\scriptsize, text=MidBlue] {$v$};
  \node[font=\scriptsize\itshape, text=DarkGrey] at (5.3,-1.1) {keeps moving};

  % Arrow 2→3
  \draw[-{Stealth}, thick, DarkGrey!60] (7.4,-2) -- (8.3,-2)
        node[midway, above, font=\scriptsize\itshape, text=DarkGrey] {force applied};

  % Ball 3: stopping
  \shade[ball color=MidBlue!80] (9.9,0.42) circle (0.42);
  \node[font=\footnotesize\bfseries, text=red!60!black] at (9.9,1.9) {Unbalanced force};
  \draw[force arrow, red!65!black] (9.48,0.42) -- (8.4,0.42)
        node[above, font=\scriptsize, text=red!65!black] {friction};
  \node[font=\scriptsize\itshape, text=DarkGrey] at (9.9,-1.1) {decelerates / stops};
\end{tikzpicture}
\end{center}}

% ── D4: Newton's Second Law — F=ma ────────────────────────────────────────────
\newcommand{\diagFma}{%
\begin{center}
\begin{tikzpicture}[font=\small]
  \node[draw=PurpleDark, fill=PurpleLight!50, rounded corners=5pt,
        minimum width=13cm, minimum height=0.9cm, align=center,
        font=\small\bfseries, text=PurpleDark]
        at (5.5,3.6)
        {Newton's Second Law: \quad $F = m \times a$ \quad $\Rightarrow$ \quad
         $a = F/m$ \quad $\Rightarrow$ \quad $m = F/a$};
  \node[draw=GoldDark, fill=GoldLight, rounded corners=5pt,
        minimum width=3.6cm, minimum height=3.2cm, align=center]
        at (0.8,1.5)
        {\textbf{\color{GoldDark}$F$\ (Force)}\\[4pt]
         \small\color{DarkGrey}Measured in newtons (N)\\
         \small\color{DarkGrey}A push or pull\\[2pt]
         \small\color{GoldDark}$F = m \times a$};
  \node[font=\large\bfseries, text=DarkGrey] at (3.3,1.5) {$=$};
  \node[draw=MidBlue, fill=LightBlue, rounded corners=5pt,
        minimum width=3.6cm, minimum height=3.2cm, align=center]
        at (5.5,1.5)
        {\textbf{\color{MidBlue}$m$\ (Mass)}\\[4pt]
         \small\color{DarkGrey}Measured in kg\\
         \small\color{DarkGrey}More mass\,=\,more\\
         \small\color{DarkGrey}inertia to overcome};
  \node[font=\large\bfseries, text=DarkGrey] at (7.75,1.5) {$\times$};
  \node[draw=TealGreen, fill=LightGreen, rounded corners=5pt,
        minimum width=3.6cm, minimum height=3.2cm, align=center]
        at (10.0,1.5)
        {\textbf{\color{TealGreen}$a$\ (Acceleration)}\\[4pt]
         \small\color{DarkGrey}Measured in m/s$^2$\\
         \small\color{DarkGrey}Bigger $F$ on same\\
         \small\color{DarkGrey}mass\,=\,bigger $a$};
\end{tikzpicture}
\end{center}}

% ── D5: Energy Types and Transformation ───────────────────────────────────────
\newcommand{\diagEnergy}{%
\begin{center}
\begin{tikzpicture}[font=\small]
  % Top row: three energy type boxes
  \node[draw=GoldDark, fill=GoldLight, rounded corners=5pt,
        minimum width=4.0cm, minimum height=2.2cm, align=center]
        at (0.0,3.2)
        {\textbf{\color{GoldDark}Kinetic Energy (KE)}\\[3pt]
         \small\color{DarkGrey}Energy of motion\\[2pt]
         \small\color{GoldDark}$KE = \tfrac{1}{2}mv^2$};
  \node[draw=MidBlue, fill=LightBlue, rounded corners=5pt,
        minimum width=4.0cm, minimum height=2.2cm, align=center]
        at (4.8,3.2)
        {\textbf{\color{MidBlue}Gravitational PE}\\[3pt]
         \small\color{DarkGrey}Stored by height\\[2pt]
         \small\color{MidBlue}$GPE = mgh$};
  \node[draw=TealGreen, fill=LightGreen, rounded corners=5pt,
        minimum width=4.0cm, minimum height=2.2cm, align=center]
        at (9.6,3.2)
        {\textbf{\color{TealGreen}Elastic PE}\\[3pt]
         \small\color{DarkGrey}Stored by stretching\\[2pt]
         \small\color{DarkGrey}springs, rubber bands};
  % Divider label
  \node[font=\small\bfseries, text=DarkGrey] at (4.8,1.85)
        {Energy Transformation in a Trebuchet};
  % Chain: GPE box
  \node[draw=OrangeDark, fill=OrangeLight, rounded corners=4pt,
        minimum width=3.0cm, minimum height=1.1cm, align=center,
        font=\small\bfseries, text=OrangeDark] (n1) at (0.8,0.7)
        {Gravitational PE\\[-1pt]{\scriptsize\color{DarkGrey}(counterweight up)}};
  % Chain: KE1 box
  \node[draw=GoldDark, fill=GoldLight, rounded corners=4pt,
        minimum width=3.0cm, minimum height=1.1cm, align=center,
        font=\small\bfseries, text=GoldDark] (n2) at (4.8,0.7)
        {Kinetic Energy\\[-1pt]{\scriptsize\color{DarkGrey}(arm swings)}};
  % Chain: KE2 box
  \node[draw=TealGreen, fill=LightGreen, rounded corners=4pt,
        minimum width=3.0cm, minimum height=1.1cm, align=center,
        font=\small\bfseries, text=TealGreen] (n3) at (8.8,0.7)
        {Kinetic Energy\\[-1pt]{\scriptsize\color{DarkGrey}(ball launches)}};
  % Waste box
  \node[draw=DarkGrey, fill=LightGrey, rounded corners=4pt,
        minimum width=2.8cm, minimum height=0.9cm, align=center,
        font=\small, text=DarkGrey] (nw) at (8.8,-0.75)
        {$+$ Thermal \& Sound\\[-1pt]{\scriptsize(wasted)}};
  % Arrows
  \draw[-{Stealth}, thick, DarkGrey] (n1.east) -- (n2.west);
  \draw[-{Stealth}, thick, DarkGrey] (n2.east) -- (n3.west);
  \draw[-{Stealth}, thick, DarkGrey] (n3.south) -- (nw.north);
\end{tikzpicture}
\end{center}}

% ── D6: Three Classes of Lever ────────────────────────────────────────────────
\newcommand{\diagLevers}{%
\begin{center}
\begin{tikzpicture}[font=\small]
  % ── Class 1: Fulcrum between E and L ──
  \node[anchor=west, font=\small\bfseries, text=NavyBlue] at (-0.3,4.85) {Class 1: Fulcrum between E and L};
  \node[anchor=east, font=\small\itshape, text=DarkGrey]  at (12.0,4.85) {e.g.\ seesaw};
  \draw[line width=4pt, DarkGrey!60, line cap=round] (1.0,4.2) -- (10.5,4.2);
  \fill[red!65!black] (5.5,4.2) -- ++(-.22,-.52) -- ++(.44,0) -- cycle;
  \node[font=\scriptsize, text=red!65!black, below] at (5.5,3.68) {F};
  \draw[force arrow, MidBlue]   (1.5,5.05) -- (1.5,4.2);
  \draw[force arrow, TealGreen] (9.5,5.05) -- (9.5,4.2);
  \node[font=\small\bfseries, text=MidBlue,   above] at (1.5,5.05) {E};
  \node[font=\small\bfseries, text=TealGreen, above] at (9.5,5.05) {L};

  % ── Class 2: Load between F and E ──
  \node[anchor=west, font=\small\bfseries, text=NavyBlue] at (-0.3,2.55) {Class 2: Load between F and E};
  \node[anchor=east, font=\small\itshape, text=DarkGrey]  at (12.0,2.55) {e.g.\ wheelbarrow};
  \draw[line width=4pt, DarkGrey!60, line cap=round] (1.0,1.9) -- (10.5,1.9);
  \fill[red!65!black] (1.0,1.9) -- ++(-.22,-.52) -- ++(.44,0) -- cycle;
  \node[font=\scriptsize, text=red!65!black, below] at (1.0,1.38) {F};
  \draw[force arrow, MidBlue]   (9.5,2.75) -- (9.5,1.9);
  \draw[force arrow, TealGreen] (5.0,2.75) -- (5.0,1.9);
  \node[font=\small\bfseries, text=MidBlue,   above] at (9.5,2.75) {E};
  \node[font=\small\bfseries, text=TealGreen, above] at (5.0,2.75) {L};

  % ── Class 3: Effort between F and L ──
  \node[anchor=west, font=\small\bfseries, text=NavyBlue] at (-0.3,0.25) {Class 3: Effort between F and L};
  \node[anchor=east, font=\small\itshape, text=DarkGrey]  at (12.0,0.25) {e.g.\ tweezers};
  \draw[line width=4pt, DarkGrey!60, line cap=round] (1.0,-0.4) -- (10.5,-0.4);
  \fill[red!65!black] (1.0,-0.4) -- ++(-.22,-.52) -- ++(.44,0) -- cycle;
  \node[font=\scriptsize, text=red!65!black, below] at (1.0,-0.92) {F};
  \draw[force arrow, MidBlue]   (5.0,0.45) -- (5.0,-0.4);
  \draw[force arrow, TealGreen] (9.5,0.45) -- (9.5,-0.4);
  \node[font=\small\bfseries, text=MidBlue,   above] at (5.0,0.45) {E};
  \node[font=\small\bfseries, text=TealGreen, above] at (9.5,0.45) {L};
\end{tikzpicture}
\end{center}}

% ── D7: Trebuchet as Class 1 Lever ────────────────────────────────────────────
\newcommand{\diagTrebuchetLever}{%
\begin{center}
\begin{tikzpicture}[font=\small]
  % Ground
  \fill[ground hatch] (0,-0.55) rectangle (12.5,0);
  \draw[thick, DarkGrey] (0,0) -- (12.5,0);
  % Uprights
  \draw[line width=3pt, DarkGrey!70] (2.8,0) -- (2.8,1.8);
  \draw[line width=3pt, DarkGrey!70] (4.4,0) -- (4.4,1.8);
  \draw[line width=3pt, DarkGrey!70] (7.6,0) -- (7.6,1.8);
  \draw[line width=3pt, DarkGrey!70] (9.2,0) -- (9.2,1.8);
  % Arm (tilted left-low to right-high)
  \draw[line width=6pt, brown!55!black, line cap=round] (2.4,2.6) -- (10.5,5.0);
  % Pivot
  \fill[DarkGrey!80] (6.2,3.6) circle (0.18);
  \node[font=\scriptsize\bfseries, text=DarkGrey, below] at (6.2,3.42) {Pivot};
  % Counterweight
  \node[draw=red!60!black, fill=red!12, rounded corners=3pt,
        minimum width=1.5cm, minimum height=1.5cm,
        font=\small\bfseries, text=red!60!black, align=center]
        at (2.0,1.4) {Heavy\\CW\\(Effort)};
  \draw[force arrow, red!65!black] (2.0,0.65) -- (2.0,-0.1);
  \node[font=\scriptsize, text=red!65!black, right] at (2.1,-0.05) {$W_1$};
  % Short arm brace
  \draw[dim arrow] (2.4,0.5) -- (6.2,0.5)
        node[midway, below, font=\scriptsize, text=red!60!black] {Short arm};
  % Ball
  \shade[ball color=MidBlue!85] (10.9,5.2) circle (0.38);
  \node[font=\scriptsize\bfseries, text=MidBlue, above right] at (11.1,5.35) {Ball (Load)};
  \draw[force arrow, red!65!black] (10.9,4.82) -- (10.9,3.8);
  \node[font=\scriptsize, text=red!65!black, right] at (11.0,4.3) {$W_2$};
  % Long arm brace
  \draw[dim arrow] (6.2,0.5) -- (10.5,0.5)
        node[midway, below, font=\scriptsize, text=MidBlue] {Long arm};
  % MA formula box
  \node[draw=GoldDark, fill=GoldLight, rounded corners=3pt,
        font=\small, text=GoldDark, align=center, inner sep=4pt]
        at (6.2,-1.3)
        {$\text{MA} = \text{Long arm} \div \text{Short arm}$};
\end{tikzpicture}
\end{center}}

% ── D8: Projectile Motion ─────────────────────────────────────────────────────
\newcommand{\diagProjectile}{%
\begin{center}
\begin{tikzpicture}[font=\small]
  % Ground
  \fill[ground hatch] (0,-0.45) rectangle (11.5,0);
  \draw[thick, DarkGrey] (0,0) -- (11.5,0);
  % Launch platform
  \fill[DarkGrey!35] (0,0) rectangle (1.3,3.0);
  \draw[thick, DarkGrey] (1.3,3.0) -- (0,3.0);
  % Height brace
  \draw[dim arrow] (-0.4,0) -- (-0.4,3.0)
        node[midway, left, font=\small\bfseries, text=red!65!black] {$h$};
  % Parabolic trajectory
  \draw[thick, MidBlue, dashed]
        (1.3,3.0) .. controls (4.2,3.8) and (8.2,2.6) .. (10.8,0);
  % Ball at launch
  \shade[ball color=MidBlue!85] (1.3,3.0) circle (0.32);
  % Horizontal velocity
  \draw[force arrow, MidBlue] (1.62,3.0) -- (3.4,3.0)
        node[above, font=\small\bfseries, text=MidBlue] {$v$ (horizontal)};
  % Weight at launch
  \draw[force arrow, red!65!black] (1.3,2.68) -- (1.3,1.4)
        node[right, font=\small\bfseries, text=red!65!black] {$W$};
  % Ball mid-trajectory
  \shade[ball color=MidBlue!60] (6.0,3.5) circle (0.22);
  \draw[-{Stealth}, thin, MidBlue]    (6.22,3.5)  -- ++(0.9,0);
  \draw[-{Stealth}, thin, red!65!black] (6.0,3.28) -- ++(0,-0.8);
  % Ball at landing
  \shade[ball color=MidBlue!85] (10.8,0.32) circle (0.32);
  % Range brace
  \draw[dim arrow] (1.3,-0.75) -- (10.8,-0.75)
        node[midway, below, font=\small\bfseries, text=TealGreen] {$R$ (horizontal range)};
  % Equations box
  \node[draw=GoldDark, fill=GoldLight, rounded corners=4pt,
        font=\small, text=GoldDark, align=center, inner sep=6pt]
        at (7.8,2.6)
        {\textbf{Key equations}\\[3pt]
         \color{NavyBlue}$h = \tfrac{1}{2}g t^2$\\[1pt]
         \color{NavyBlue}$g = 2h \div t^2$\\[1pt]
         \color{NavyBlue}$R = v \times t$\\[2pt]
         \color{DarkGrey}\scriptsize$g = 9.8$\,m/s$^2$};
\end{tikzpicture}
\end{center}}

% ── Experiment 2: Trolley and pulley setup ─────────────────────────────────────
\newcommand{\diagExpTwoSetup}{%
\begin{center}
\begin{tikzpicture}[font=\small, x=1cm, y=1cm]
  % bench ends just before the pulley, so the hanging mass falls past its edge
  \fill[ground hatch] (0,-0.4) rectangle (8.95,0);
  \draw[thick, DarkGrey] (0,0) -- (8.95,0) -- (8.95,-0.4);
  \fill[DarkGrey!15] (0.3,0) rectangle (8.6,0.18);
  \draw[DarkGrey!60] (0.3,0.18) -- (8.6,0.18);
  \node[draw=NavyBlue, fill=LightBlue, rounded corners=3pt,
        minimum width=2.4cm, minimum height=1.0cm,
        font=\small\bfseries, text=NavyBlue] (trolley) at (3.5,0.68) {Trolley ($M$)};
  \fill[DarkGrey!60] (2.55,0.18) circle (0.13);
  \fill[DarkGrey!60] (4.45,0.18) circle (0.13);
  \draw[dashed, DarkGrey!50, thick] (2.3,0.68) -- (0.3,0.68);
  \node[font=\scriptsize\itshape, text=DarkGrey, anchor=east] at (0.22,0.68) {ticker tape};
  % pulley clamped to the bench edge: string runs over its top, down its right side
  \draw[thick, DarkGrey] (8.95,0.0) -- (8.95,0.34);
  \draw[thick, DarkGrey!70] (trolley.east) -- (8.95,0.68);
  \fill[DarkGrey!25] (8.95,0.34) circle (0.34);
  \draw[thick, DarkGrey] (8.95,0.34) circle (0.34);
  \draw[thick, DarkGrey] (8.95,0.34) circle (0.11);
  \node[font=\scriptsize, text=DarkGrey, above right] at (9.15,0.6) {pulley};
  \draw[thick, DarkGrey!70] (9.29,0.34) -- (9.29,-1.77);
  \node[draw=OrangeDark, fill=OrangeLight, rounded corners=3pt,
        minimum width=1.3cm, minimum height=0.85cm,
        font=\small\bfseries, text=OrangeDark] (hm) at (9.29,-2.2) {$m$};
  \node[font=\scriptsize, text=OrangeDark, anchor=west] at (9.99,-2.2) {hanging mass};
  \draw[force arrow, red!65!black] (hm.south) -- ++(0,-1.0)
        node[right, font=\scriptsize\bfseries, text=red!65!black] {$W\!=\!mg$};
  \draw[force arrow, TealGreen] (9.29,-1.6) -- (9.29,-0.9)
        node[right, font=\scriptsize\bfseries, text=TealGreen] {$T$};
  \draw[force arrow, OrangeDark] (trolley.east) ++(0.1,0) -- ++(1.5,0)
        node[above, font=\scriptsize\bfseries, text=OrangeDark] {$T$};
  \draw[force arrow, PurpleDark] (trolley.north) ++(0.0,0.1) -- ++(1.8,0)
        node[above, font=\scriptsize\bfseries, text=PurpleDark] {$a$};
  \node[font=\scriptsize\itshape, text=DarkGrey] at (3.5,2.05)
        {total mass $M + m$ kept constant (masses move from trolley to hanger)};
  \node[font=\scriptsize\itshape, text=DarkGrey, anchor=west] at (0.3,-0.65) {bench / runway};
\end{tikzpicture}
\end{center}}

% ── Experiment 2: FBD of trolley and hanging mass side by side ────────────────
\newcommand{\diagExpTwoFBD}{%
\begin{center}
\begin{tikzpicture}[font=\small]
  \node[font=\small\bfseries, text=NavyBlue] at (0,3.5) {FBD: Trolley ($M$)};
  \node[draw=NavyBlue, fill=LightBlue, rounded corners=4pt,
        minimum width=2.6cm, minimum height=1.2cm,
        font=\small\bfseries, text=NavyBlue] (T) at (0,0) {Trolley ($M$)};
  \draw[force arrow, TealGreen] (T.north) -- ++(0,2.0)
        node[right, font=\small\bfseries, text=TealGreen] {$N$ (normal)};
  \draw[force arrow, red!65!black] (T.south) -- ++(0,-2.0)
        node[right, font=\small\bfseries, text=red!65!black] {$W\!=\!Mg$};
  \draw[force arrow, OrangeDark] (T.east) -- ++(2.8,0)
        node[above, font=\small\bfseries, text=OrangeDark] {$T$ (tension)};
  \draw[thin arrow, DarkGrey!55, dashed] (T.west) -- ++(-1.8,0)
        node[above, font=\scriptsize\itshape, text=DarkGrey!65] {friction $\approx 0$};
  \node[font=\scriptsize\itshape, text=DarkGrey, anchor=east] at (-0.2,1.5)
        {$N = Mg$ (vertical balance)};
  \node[draw=OrangeDark, fill=OrangeLight, rounded corners=3pt,
        font=\small, text=OrangeDark, inner sep=3pt]
        at (0,-3.6) {$T = Ma$};

  \node[font=\small\bfseries, text=NavyBlue] at (8.5,3.5) {FBD: Hanging mass ($m$)};
  \node[draw=OrangeDark, fill=OrangeLight, rounded corners=4pt,
        minimum width=2.4cm, minimum height=1.1cm,
        font=\small\bfseries, text=OrangeDark] (H) at (8.5,0) {Mass ($m$)};
  \draw[force arrow, TealGreen] (H.north) -- ++(0,1.6)
        node[right, font=\small\bfseries, text=TealGreen] {$T$ (tension)};
  \draw[force arrow, red!65!black] (H.south) -- ++(0,-2.6)
        node[right, font=\small\bfseries, text=red!65!black] {$W\!=\!mg$};
  \node[font=\scriptsize\itshape, text=red!60!black] at (8.5,2.9)
        {$mg > T$ --- forces unbalanced};
  \draw[force arrow, red!65!black, line width=1.2pt, opacity=0.5]
        (H.south) ++(0.45,-0.3) -- ++(0,-0.9)
        node[right, font=\scriptsize\bfseries, text=red!65!black] {net $F$};
  \node[draw=red!55!black, fill=red!8, rounded corners=3pt,
        font=\small, text=red!55!black, inner sep=3pt, align=center]
        at (8.5,-4.0) {$mg - T = ma$};
\end{tikzpicture}
\end{center}}

% ── Experiment 2: Comparing force sizes (50g / 100g / 200g) ──────────────────
\newcommand{\diagExpTwoForces}{%
\begin{center}
\begin{tikzpicture}[font=\small]
  \node[font=\small\bfseries, text=NavyBlue] at (7.0,1.4)
        {Effect of increasing hanging mass --- arrow length shows force size};
  \node[font=\small\bfseries, text=GoldDark] at (1.5,0.5) {$m = 50$\,g};
  \node[draw=NavyBlue, fill=LightBlue, rounded corners=3pt,
        minimum width=2.0cm, minimum height=0.9cm,
        font=\small\bfseries, text=NavyBlue] (t1) at (1.5,-0.8) {$M$};
  \draw[force arrow, OrangeDark] (t1.east) -- ++(1.0,0)
        node[above, font=\scriptsize, text=OrangeDark] {$T_1$};
  \draw[force arrow, PurpleDark] (t1.south east) -- ++(0.7,0)
        node[below, font=\scriptsize, text=PurpleDark] {$a_1$};
  \node[draw=OrangeDark, fill=OrangeLight, rounded corners=2pt,
        minimum width=0.9cm, minimum height=0.7cm,
        font=\scriptsize\bfseries, text=OrangeDark] at (3.0,-2.5) {50\,g};
  \draw[force arrow, red!65!black] (3.0,-2.85) -- ++(0,-0.7)
        node[right, font=\scriptsize, text=red!65!black] {$W_1$};

  \node[font=\small\bfseries, text=GoldDark] at (6.0,0.5) {$m = 100$\,g};
  \node[draw=NavyBlue, fill=LightBlue, rounded corners=3pt,
        minimum width=2.0cm, minimum height=0.9cm,
        font=\small\bfseries, text=NavyBlue] (t2) at (6.0,-0.8) {$M$};
  \draw[force arrow, OrangeDark] (t2.east) -- ++(1.7,0)
        node[above, font=\scriptsize, text=OrangeDark] {$T_2$};
  \draw[force arrow, PurpleDark] (t2.south east) -- ++(1.2,0)
        node[below, font=\scriptsize, text=PurpleDark] {$a_2$};
  \node[draw=OrangeDark, fill=OrangeLight, rounded corners=2pt,
        minimum width=1.0cm, minimum height=0.85cm,
        font=\scriptsize\bfseries, text=OrangeDark] at (8.1,-2.5) {100\,g};
  \draw[force arrow, red!65!black] (8.1,-2.925) -- ++(0,-1.1)
        node[right, font=\scriptsize, text=red!65!black] {$W_2$};

  \node[font=\small\bfseries, text=GoldDark] at (10.8,0.5) {$m = 200$\,g};
  \node[draw=NavyBlue, fill=LightBlue, rounded corners=3pt,
        minimum width=2.0cm, minimum height=0.9cm,
        font=\small\bfseries, text=NavyBlue] (t3) at (10.8,-0.8) {$M$};
  \draw[force arrow, OrangeDark] (t3.east) -- ++(2.4,0)
        node[above, font=\scriptsize, text=OrangeDark] {$T_3$};
  \draw[force arrow, PurpleDark] (t3.south east) -- ++(1.8,0)
        node[below, font=\scriptsize, text=PurpleDark] {$a_3$ (largest)};
  \node[draw=OrangeDark, fill=OrangeLight, rounded corners=2pt,
        minimum width=1.1cm, minimum height=1.1cm,
        font=\scriptsize\bfseries, text=OrangeDark] at (13.6,-2.5) {200\,g};
  \draw[force arrow, red!65!black] (13.6,-3.05) -- ++(0,-1.6)
        node[right, font=\scriptsize, text=red!65!black] {$W_3$};

  \node[draw=TealGreen, fill=LightGreen, rounded corners=4pt,
        text width=13.4cm, inner sep=5pt,
        font=\small, text=TealGreen, align=center]
        at (7.0,-5.95)
        {Prediction: when the hanging mass doubles from 50\,g to 100\,g, the tension doubles,
         so the acceleration should double too. When the mass stays the same,
         force and acceleration are directly proportional.};
\end{tikzpicture}
\end{center}}

% ── Experiment 2: Combined system equation derivation ─────────────────────────
\newcommand{\diagExpTwoSystem}{%
\begin{center}
\begin{tikzpicture}[font=\small]
  \node[draw=NavyBlue, fill=LightBlue!30, rounded corners=5pt,
        minimum width=13cm, minimum height=6.4cm] at (5.5,-0.4) {};
  \node[font=\small\bfseries, text=NavyBlue] at (5.5,2.4)
        {System analysis: combining both free body diagrams};
  \node[font=\small\bfseries, text=NavyBlue, anchor=north west] at (0.4,1.8) {Trolley ($M$):};
  \node[font=\small, text=DarkGrey, anchor=north west] at (0.4,1.3) {Horizontal net force $= T$};
  \node[font=\small, text=NavyBlue, anchor=north west] at (0.4,0.8)
        {$\Rightarrow\quad T = Ma \quad\cdots(1)$};
  \node[font=\small\bfseries, text=OrangeDark, anchor=north west] at (6.2,1.8) {Hanging mass ($m$):};
  \node[font=\small, text=DarkGrey, anchor=north west] at (6.2,1.3) {Net force $= mg - T$ downward};
  \node[font=\small, text=OrangeDark, anchor=north west] at (6.2,0.8)
        {$\Rightarrow\quad mg - T = ma \quad\cdots(2)$};
  \node[font=\small\bfseries, text=TealGreen] at (5.5,-0.2)
        {Add (1) and (2) to eliminate $T$:};
  \node[font=\small, text=NavyBlue] at (5.5,-0.75)
        {$mg = Ma + ma = (M+m)\,a$};
  \node[font=\small\bfseries, text=NavyBlue] at (5.5,-1.6)
        {$\displaystyle a = \frac{mg}{M+m}$};
  \node[draw=GoldDark, fill=GoldLight, rounded corners=3pt,
        font=\small, text=GoldDark, inner sep=5pt, align=center]
        at (5.5,-2.85)
        {When $M \gg m$, the total mass $\approx M$, so
         $a \approx \dfrac{mg}{M} = \dfrac{F}{M}$ --- confirming $F = Ma$};
\end{tikzpicture}
\end{center}}


% ── New in this update: spring-scale and force-arrow scale diagrams ──────────
\newcommand{\diagSpringScales}{%
\begin{center}
\begin{tikzpicture}[font=\footnotesize, x=1cm, y=1cm]
  % --- Scale A: 0-10 N, pointer at 4 N ---
  \draw[line width=1pt, Ink] (0,0) rectangle (1.5,6);
  \node[font=\small\bfseries, text=Ink] at (0.75,6.35) {A};
  \draw[line width=1pt, Ink] (0.5,6) -- (0.5,6.5);
  \draw[line width=1pt, Ink] (0.75,-0.5) circle (0.22);
  \draw[line width=1pt, Ink] (0.75,0) -- (0.75,-0.28);
  \foreach \v in {0,1,2,3,4,5,6,7,8,9,10}{
    \pgfmathsetmacro{\yy}{5.5-\v*0.5}
    \draw[line width=0.7pt, Ink] (0.15,\yy) -- (0.55,\yy);
    \node[anchor=west, font=\tiny, text=DarkGrey] at (0.62,\yy) {\v};
  }
  \draw[line width=2pt, red!70!black] (0.1,3.5) -- (1.4,3.5);
  \node[anchor=west, font=\tiny\bfseries, text=red!70!black] at (1.5,3.5) {};
  \node[anchor=north, font=\tiny, text=DarkGrey] at (0.75,-0.85) {0--10\,N};

  % --- Scale B: 0-5 N, pointer at 2.5 N ---
  \begin{scope}[xshift=4cm]
    \draw[line width=1pt, Ink] (0,0) rectangle (1.5,6);
    \node[font=\small\bfseries, text=Ink] at (0.75,6.35) {B};
    \draw[line width=1pt, Ink] (0.5,6) -- (0.5,6.5);
    \draw[line width=1pt, Ink] (0.75,-0.5) circle (0.22);
    \draw[line width=1pt, Ink] (0.75,0) -- (0.75,-0.28);
    \foreach \v in {0,1,2,3,4,5}{
      \pgfmathsetmacro{\yy}{5.5-\v*1.0}
      \draw[line width=0.7pt, Ink] (0.15,\yy) -- (0.55,\yy);
      \node[anchor=west, font=\tiny, text=DarkGrey] at (0.62,\yy) {\v};
    }
    \foreach \v in {0.5,1.5,2.5,3.5,4.5}{
      \pgfmathsetmacro{\yy}{5.5-\v*1.0}
      \draw[line width=0.5pt, DarkGrey] (0.25,\yy) -- (0.45,\yy);
    }
    \draw[line width=2pt, red!70!black] (0.1,3.0) -- (1.4,3.0);
    \node[anchor=north, font=\tiny, text=DarkGrey] at (0.75,-0.85) {0--5\,N};
  \end{scope}

  % --- Scale C: 0-20 N, pointer at 13 N ---
  \begin{scope}[xshift=8cm]
    \draw[line width=1pt, Ink] (0,0) rectangle (1.5,6);
    \node[font=\small\bfseries, text=Ink] at (0.75,6.35) {C};
    \draw[line width=1pt, Ink] (0.5,6) -- (0.5,6.5);
    \draw[line width=1pt, Ink] (0.75,-0.5) circle (0.22);
    \draw[line width=1pt, Ink] (0.75,0) -- (0.75,-0.28);
    \foreach \v in {0,2,4,6,8,10,12,14,16,18,20}{
      \pgfmathsetmacro{\yy}{5.5-\v*0.25}
      \draw[line width=0.7pt, Ink] (0.15,\yy) -- (0.55,\yy);
      \node[anchor=west, font=\tiny, text=DarkGrey] at (0.62,\yy) {\v};
    }
    \draw[line width=2pt, red!70!black] (0.1,2.25) -- (1.4,2.25);
    \node[anchor=north, font=\tiny, text=DarkGrey] at (0.75,-0.85) {0--20\,N};
  \end{scope}

  % Instruction note
  \node[anchor=west, font=\footnotesize\itshape, text=DarkGrey, align=left]
        at (10.4,3) {Read the value where\\the red pointer sits.\\Check what one small\\division is worth first.};
\end{tikzpicture}
\end{center}}

\newcommand{\diagArrowScale}{%
\begin{center}
\begin{tikzpicture}[font=\footnotesize, x=1cm, y=1cm]
  % Scale key
  \node[anchor=west, font=\small\bfseries, text=Ink] at (0,3.4) {Scale: 1\,cm represents 5\,N};
  \draw[line width=1.2pt, Ink] (0,3.0) -- (1,3.0);
  \draw[line width=0.8pt, Ink] (0,2.88) -- (0,3.12);
  \draw[line width=0.8pt, Ink] (1,2.88) -- (1,3.12);
  \node[anchor=west, font=\footnotesize, text=DarkGrey] at (1.2,3.0) {= 1\,cm = 5\,N};

  % 5 N arrow
  \draw[force arrow, MidBlue] (0,2.0) -- (1,2.0);
  \node[anchor=west, font=\footnotesize\bfseries, text=MidBlue] at (1.3,2.0) {5\,N \; (1\,cm long)};

  % 10 N arrow
  \draw[force arrow, TealGreen] (0,1.2) -- (2,1.2);
  \node[anchor=west, font=\footnotesize\bfseries, text=TealGreen] at (2.3,1.2) {10\,N \; (2\,cm long)};

  % 15 N arrow
  \draw[force arrow, OrangeDark] (0,0.4) -- (3,0.4);
  \node[anchor=west, font=\footnotesize\bfseries, text=OrangeDark] at (3.3,0.4) {15\,N \; (3\,cm long)};

  % Worked box on right
  \node[draw=GoldDark, fill=GoldLight, rounded corners=3pt, inner sep=5pt,
        font=\footnotesize, text=GoldDark, align=left, anchor=north west]
        at (8.2,3.3)
        {\textbf{To find the length:}\\[2pt]
         length $=$ force $\div$ scale\\[2pt]
         e.g.\ $15 \div 5 = 3$\,cm};
\end{tikzpicture}
\end{center}}


% ── Diagrams added with Lessons 4–6 (Newton's laws) ───────────────────────
% ── Newton's Third Law: two skaters, and a rocket with its exhaust ────────────
\newcommand{\diagThirdLaw}{%
\begin{center}
\begin{tikzpicture}[font=\small]
  %% ── LEFT: two skaters push apart ──
  \node[font=\small\bfseries, text=NavyBlue] at (2.7,2.95) {Two skaters push apart};
  \fill[LightBlue] (0.3,-0.3) rectangle (5.1,0);
  \draw[thick, Sky] (0.3,0) -- (5.1,0);
  \node[font=\scriptsize\itshape, text=DarkGrey] at (2.7,-0.55) {ice};
  \node[draw=NavyBlue, fill=LightBlue, rounded corners=4pt,
        minimum width=1.4cm, minimum height=2.1cm,
        font=\small\bfseries, text=NavyBlue] (A) at (2.0,1.1) {A};
  \node[draw=TealGreen, fill=LightGreen, rounded corners=4pt,
        minimum width=1.4cm, minimum height=2.1cm,
        font=\small\bfseries, text=TealGreen] (B) at (3.4,1.1) {B};
  \draw[force arrow, OrangeDark] (A.west) -- ++(-1.3,0)
        node[above, font=\scriptsize\bfseries, text=OrangeDark, align=center] {B pushes A\\60\,N};
  \draw[force arrow, OrangeDark] (B.east) -- ++(1.3,0)
        node[above, font=\scriptsize\bfseries, text=OrangeDark, align=center] {A pushes B\\60\,N};

  %% ── RIGHT: a rocket in space and its exhaust gas ──
  \node[font=\small\bfseries, text=NavyBlue] at (11.0,2.95) {A rocket and its exhaust gas};
  % exhaust gas (drawn first)
  \shade[inner color=Amber!80, outer color=OrangeLight] (9.3,1.1) ellipse (0.6 and 0.42);
  \node[font=\scriptsize\bfseries, text=GoldDark] at (9.3,1.1) {gas};
  % fins
  \fill[NavyBlue!75] (9.95,1.45) -- (9.7,1.85) -- (10.35,1.45) -- cycle;
  \fill[NavyBlue!75] (9.95,0.75) -- (9.7,0.35) -- (10.35,0.75) -- cycle;
  % body and nose
  \draw[NavyBlue, fill=LightBlue, thick] (9.95,0.75) rectangle (11.75,1.45);
  \fill[NavyBlue] (11.75,0.75) -- (12.25,1.1) -- (11.75,1.45) -- cycle;
  \node[font=\scriptsize\bfseries, text=NavyBlue] at (10.9,1.1) {rocket};
  % the pair
  \draw[force arrow, OrangeDark] (8.7,1.1) -- ++(-1.3,0)
        node[above, font=\scriptsize\bfseries, text=OrangeDark, align=center] {rocket pushes\\gas back};
  \draw[force arrow, OrangeDark] (12.25,1.1) -- ++(1.3,0)
        node[above, font=\scriptsize\bfseries, text=OrangeDark, align=center] {gas pushes\\rocket forward};
  \node[font=\scriptsize\itshape, text=DarkGrey] at (10.8,-0.55) {in space --- no air needed};
\end{tikzpicture}
\end{center}}

% ── Newton's Second Law: formula triangle ──────────────────────────────────────
\newcommand{\diagFmaTriangle}{%
\begin{center}
\begin{tikzpicture}[font=\small]
  \draw[thick, GoldDark, fill=GoldLight, line join=round] (0,0) -- (3.8,0) -- (1.9,3.0) -- cycle;
  \draw[thick, GoldDark] (0.76,1.2) -- (3.04,1.2);
  \draw[thick, GoldDark] (1.9,0) -- (1.9,1.2);
  \node[text=GoldDark]  at (1.9,1.85) {\LARGE $\boldsymbol{F}$};
  \node[text=MidBlue]   at (1.22,0.52) {\LARGE $\boldsymbol{m}$};
  \node[text=TealGreen] at (2.58,0.52) {\LARGE $\boldsymbol{a}$};
  \node[text=DarkGrey, fill=GoldLight, inner sep=1pt] at (1.9,0.52) {$\times$};
  % how to use it
  \node[anchor=west, align=left, text=Ink, font=\normalsize] at (4.8,1.5)
       {\textbf{Cover} $\boldsymbol{F}$:\quad $F = m \times a$\\[7pt]
        \textbf{Cover} $\boldsymbol{a}$:\quad $a = F \div m$\\[7pt]
        \textbf{Cover} $\boldsymbol{m}$:\quad $m = F \div a$};
  \node[anchor=west, align=left, text=DarkGrey, font=\footnotesize\itshape] at (4.8,-0.1)
       {Side by side means multiply. One above the other means divide.};
\end{tikzpicture}
\end{center}}


% ══════════════════════════════════════════════════════════════════════════════
%  LESSON STRUCTURE (lesson booklet, Oct 2026)
%  Every lesson runs in the same four parts, in this order:
%    \lessonpart{A}{Reading and comprehension}{Teal}   \label{read:N}
%    \lessonpart{B}{Theory and examples}{Blue}         \label{theory:N}
%    \lessonpart{C}{Practice}{Amber}                    \label{prac:N}
%    \lessonpart{D}{Above level}{Violet}                \label{ext:N}
%  Above-level / higher-order questions live ONLY in Part D, at the end.
% ══════════════════════════════════════════════════════════════════════════════
\newcommand{\lessonpart}[3]{%
  \par\if#1A\needspace{10\baselineskip}\vspace{10pt}\else
    \if#1C\notesfill\fi\clearpage\fi
  \begin{tcolorbox}[enhanced, breakable=false, arc=2mm, boxrule=0pt,
    colback=#3!14, borderline west={4pt}{0pt}{#3},
    left=4mm, right=3mm, top=1.6mm, bottom=1.6mm]
  {\headingfont\bfseries\color{#3!60!black}\large Part #1 \textbullet{} #2}
  \end{tcolorbox}\vspace{2pt}}

% Curriculum codes under a lesson banner
\newcommand{\lessoncodes}[1]{%
  {\small\headingfont\color{DarkGrey}Victorian Curriculum 2.0:\ \textbf{#1}}\par\vspace{4pt}}

% Two short tasks side by side, each with its own working box
% \taskpair{text}{box height}{text}{box height}
\newcommand{\taskpair}[4]{%
  \par\needspace{9\baselineskip}\vspace{4pt}\noindent
  \edef\abgrp{\number\numexpr\value{abox}+1\relax}%
  \begin{minipage}[t]{0.485\linewidth}
    \task #1\par\vspace{2pt}\answerbox{#2}
  \end{minipage}\hfill
  \begin{minipage}[t]{0.485\linewidth}
    \task #3\par\vspace{2pt}\answerbox{#4}
  \end{minipage}\par\gdef\abgrp{0}\vspace{4pt}}

% Answer boxes never split across a page (a half box on its own page is useless)
% ── Page filling ────────────────────────────────────────────────────────────
% Every writing box gets an id. Its page is logged to \jobname.abp; fill.py
% measures the empty space left at the foot of each page and shares it out
% among that page's boxes via \jobname.abx (\abset{id}{extra}). Side-by-side
% boxes share a group so they grow together.
\newcounter{abox}
\newwrite\abpos \immediate\openout\abpos=\jobname.abp
\newcommand{\abset}[2]{\expandafter\gdef\csname abx@#1\endcsname{#2}}
\InputIfFileExists{\jobname.abx}{}{}
\def\abgrp{0}
\newcommand{\abnext}{\stepcounter{abox}%
  \edef\abcur{\ifcsname abx@\theabox\endcsname\csname abx@\theabox\endcsname\else 0pt\fi}}
\newcommand{\ablog}{\edef\abtmp{\noexpand\write\abpos{\theabox\space\abgrp\space\noexpand\thepage}}\abtmp}
\newcommand{\abh}[1]{\dimexpr#1+\abcur\relax}
\renewcommand{\answerbox}[1]{%
  \abnext
  \begin{tcolorbox}[enhanced, breakable=false, arc=2mm, colback=Paper, colframe=Line,
    left=2mm, right=2mm, top=1mm, bottom=1mm, boxrule=1pt, height=\abh{#1}]
  \end{tcolorbox}\ablog%
}
% \drawbox{height}{contents}: white box for a diagram or drawing
\newcommand{\drawbox}[2]{%
  \abnext
  \begin{tcolorbox}[colback=white, colframe=MidGrey, boxrule=0.5pt, height=\abh{#1}]#2\end{tcolorbox}\ablog}
% Notes box at the end of Part B: fills whatever is left of the page (if worth it)
\newcommand{\notesfill}{\par
  \ifdim\dimexpr\pagegoal-\pagetotal\relax>3.2cm
    \vspace*{\fill}%
    \begin{tcolorbox}[enhanced, breakable=false, arc=2mm, colback=Paper, colframe=Line,
      left=2mm, right=2mm, top=1mm, bottom=1mm, boxrule=1pt,
      height=\dimexpr\pagegoal-\pagetotal-1.0cm\relax]
      {\small\headingfont\color{DarkGrey}Notes}
    \end{tcolorbox}%
  \fi}



% Centred table: wrap every content table in \begin{ctab} ... \end{ctab}
% (figures already centre themselves inside their \diag... macros)
% ctab also registers write-in tables (S column) for page filling: the extra
% height is per row, logged with the row count.
\newenvironment{ctab}{\par\smallskip\begingroup\centering
  \abnext\let\tbcur\abcur\global\wtrows0 }{\par
  \ifnum\wtrows>0 \edef\abtmp{\noexpand\write\abpos{\theabox\space\abgrp\space\noexpand\thepage\space\the\wtrows}}\abtmp\fi
  \endgroup\smallskip}

% ── Booklet formatting rules (Oct 2026) ─────────────────────────────────────
% 1. No marks are shown anywhere in the booklet.
\renewcommand{\markscount}[1]{}
% 2. A task's heading sits on its own line; its text starts on the next
%    line, left aligned. Parts (a), (b)… each start a new line: use \pt{a}.
\RenewDocumentCommand{\task}{o}{\par\noindent\refstepcounter{tasknum}\label{task:\thetasknum\IfValueT{#1}{#1}}\textbf{Task \thetasknum\IfValueT{#1}{#1}}\par\nobreak\noindent\ignorespaces}
\renewcommand{\taskpart}[1]{\par\noindent\label{task:\thetasknum#1}\textbf{Task \thetasknum#1}\par\nobreak\noindent\ignorespaces}
\renewcommand{\exttask}{\par\noindent\refstepcounter{tasknum}\label{task:\thetasknum}\textbf{Extension \thetasknum}\par\nobreak\noindent\ignorespaces}
\newcommand{\pt}[1]{\par\noindent\textbf{(#1)}\ \ignorespaces}
% 3. Tables have vertical rules between columns and down both sides
%    (write | between and around the columns, e.g. {|L{3cm}|C{2cm}|}).
%    Keep the normal padding at the outer edges (no @{} there): colortbl paints
%    header colour one \tabcolsep past each edge, so trimming the padding makes
%    the header overhang the grid.
\arrayrulecolor{MidGrey}
\setlength{\arrayrulewidth}{0.6pt}
% 4. Write-in tables: W/Z are left/centred columns that centre text
%    vertically; put the S column first (see below).
\newcolumntype{W}[1]{>{\RaggedRight\arraybackslash}m{#1}}
\newcolumntype{Z}[1]{>{\Centering\arraybackslash}m{#1}}
\def\tbcur{0pt}
\newcount\wtrows
\newcommand{\wh}{\global\advance\wtrows1 \rule[\dimexpr-0.75cm-\tbcur/2\relax]{0pt}{\dimexpr1.5cm+\tbcur\relax}}
% S = zero-width first column that makes every row of a write-in table
% at least 1.5 cm tall with its text vertically centred. Start each data
% row with "&"; start the header row with "\hd &" so it stays slim.
\newif\ifnowh
\newcolumntype{S}{@{}>{}c<{\ifnowh\else\wh\fi\global\nowhfalse}@{\hspace{\tabcolsep}}}
\newcommand{\hd}{\global\nowhtrue}
% 5. The first row of a multi-row write-in table is a worked scaffold in red.
\definecolor{ScaffoldRed}{HTML}{C0392B}
\newcommand{\scaf}[1]{\textcolor{ScaffoldRed}{\itshape #1}}
% 6. Tables are a light grid (rules between rows and columns) on white —
%    row shading is off, because colortbl paints shading over column rules.
\renewcommand{\talt}{}
\renewcommand{\tmid}{}
\begin{document}
\setlength{\parindent}{0pt}
\setlength{\parskip}{4pt}

% ══════════════════════════════════════════════════════════════════════════════
%  COVER PAGE
% ══════════════════════════════════════════════════════════════════════════════
\thispagestyle{plain}

\begin{tcolorbox}[
  enhanced, breakable=false, arc=4mm,
  colback=NavyBlue, colframe=NavyBlue,
  left=6mm, right=6mm, top=9mm, bottom=9mm,
  boxrule=0pt
]
  \begin{center}
    {\headingfont\bfseries\color{Sky}\footnotesize\MakeUppercase{Lesson Booklet \textbullet\ Years 7 and 8}}\\[10pt]
    {\headingfont\bfseries\color{white}\fontsize{40}{44}\selectfont Forces and Motion}\\[8pt]
    {\large\color{white!85!NavyBlue}\itshape Types of force \textbullet\ Newton's laws \textbullet\ energy \textbullet\ levers \textbullet\ projectiles}\\[14pt]
    {\headingfont\color{white!72!NavyBlue}\normalsize Victorian Curriculum F--10 Version 2.0 \textbar{}
     VC2S8U13 \textbar{} VC2S8U14 \textbar{} VC2S8U15 \textbar{} VC2S8I01--I08}
  \end{center}
\end{tcolorbox}

\vspace{10pt}

% Student details
\begin{tabular}{@{}p{3.5cm} p{6cm} p{2cm} p{3.5cm}@{}}
  \textbf{Name:}    & \underline{\hspace{5.8cm}} & \textbf{Class:}   & \underline{\hspace{3.3cm}} \\[8pt]
  \textbf{Teacher:} & \underline{\hspace{5.8cm}} & \textbf{Year level:} & \underline{\hspace{3.3cm}} \\
\end{tabular}

\vspace{10pt}

\begin{infobox}{How to use this booklet}
There is one section for each lesson, and every lesson works the same way:

\medskip
\begin{tabular}{@{}ll@{}}
  \textcolor{Teal}{\rule{10pt}{10pt}}\ \textbf{Part A --- Reading} & a short passage and comprehension questions to start the lesson \\
  \textcolor{Blue}{\rule{10pt}{10pt}}\ \textbf{Part B --- Theory} & the ideas, with worked examples to follow step by step \\
  \textcolor{Amber}{\rule{10pt}{10pt}}\ \textbf{Part C --- Practice} & questions like the examples, so you can build the skill \\
  \textcolor{Violet}{\rule{10pt}{10pt}}\ \textbf{Part D --- Above level} & harder, higher-order questions --- Year 8 and extension \\
\end{tabular}

\medskip
Experiment lessons follow the same four parts: a reading about the science, planning the experiment, results and analysis, then above-level questions. Everything leads to the trebuchet investigation at the end of the unit.
\end{infobox}

\vspace{8pt}

% Contents — rows are added as each lesson is revised
\begin{ctab}
\begin{tabular}{|C{1.4cm}|L{9.6cm}|C{1.6cm}|}
  \thead{Lesson} & \thead{Title} & \thead{Page} \\ \hline
  \talt 1 & What is a Force? & p.\ \pageref{sec:1} \\ \hline
        2 & Force Diagrams and Net Force & p.\ \pageref{sec:2} \\ \hline
  \talt 3 & Experiment 1 --- Investigating Friction & p.\ \pageref{sec:3} \\ \hline
        4 & Newton's First and Third Laws & p.\ \pageref{sec:4} \\ \hline
  \talt 5 & Newton's Second Law: $F = ma$ & p.\ \pageref{sec:5} \\ \hline
        6 & Experiment 2 --- Investigating Newton's Second Law & p.\ \pageref{sec:6} \\ \hline
\end{tabular}
\end{ctab}

\newpage
% ══════════════════════════════════════════════════════════════════════════════
%  LESSON 1 — WHAT IS A FORCE?                                    (Week 1 · P1)
% ══════════════════════════════════════════════════════════════════════════════
\sectiondiv{What is a Force?}{Pushes and pulls \textbullet{} Measuring force \textbullet{} Mass and weight \textbullet{} Contact and non-contact forces}\label{sec:1}
\lessoncodes{VC2S8U14 \textbullet{} VC2S8I04}

% ─── PART A ───────────────────────────────────────────────────────────────────
\lessonpart{A}{Reading and comprehension}{Teal}\label{read:1}

\begin{readingbox}{Reading --- The Forces You Cannot See}
Right now, while you sit perfectly still, you are being pushed and pulled from several directions at once. Gravity is pulling every part of your body downwards towards the centre of the Earth. At the same time, the chair beneath you is pushing upwards on you with exactly the same strength. Because these two forces are equal and opposite, you stay put. If the chair suddenly vanished, the pushing force would disappear, gravity would win, and you would fall.

People often think that astronauts on the International Space Station float because there is no gravity in space. This is not true. At the height the Space Station orbits, Earth's gravity is still about 90\% as strong as it is on the ground. Astronauts float because they --- and the whole Space Station --- are falling around the Earth at enormous speed, and they fall together. Everything inside falls at the same rate, so nothing presses against anything else. The astronauts do not feel the push of a floor, so they seem weightless.

Forces are invisible, so scientists need a way to record them. The solution is simple but powerful: draw each force as an arrow. The direction of the arrow shows which way the force acts, and the length of the arrow shows how strong it is. A picture full of invisible pushes and pulls becomes something you can measure, compare and calculate with.
\end{readingbox}

\rcreset
\rcq{According to the passage, name the two forces acting on you while you sit still, and state the direction of each.}{1.6cm}
\rcq{The passage says that if the chair vanished, ``gravity would win''. Explain what this means, using the idea of two forces that are no longer equal.}{1.8cm}
\rcq{Explain why astronauts on the Space Station appear to float, even though gravity is still acting on them.}{2.2cm}
\rcq{Does an astronaut on the Space Station still have weight? Use a number from the passage to support your answer.}{1.6cm}
\rcq{The passage says drawing forces as arrows is ``simple but powerful''. What two pieces of information does a force arrow give you?}{1.6cm}
\rcq{Which three-word phrase in the first paragraph describes two forces that are the same size but act in opposite directions?}{1.4cm}
\rcq{Forces are invisible. Describe one way you could tell that a force is acting on an object, even though you cannot see the force itself.}{1.8cm}
\rcq{Imagine the Space Station suddenly stopped moving around the Earth. Using the passage, predict what would happen to the station and the astronauts inside, and explain why.}{2.2cm}

% ─── PART B ───────────────────────────────────────────────────────────────────
\lessonpart{B}{Theory and examples}{Blue}\label{theory:1}

\subsection{Every Force is a Push or a Pull}\label{sub:1.1}

This is the whole idea, and it really is that simple. Whenever one object affects another --- by touching it, or from a distance --- it does so by \textbf{pushing} it or \textbf{pulling} it. There is no third option.

\begin{skillbox}{Think about it}
A magnet \emph{pulls} a paperclip. A hand \emph{pushes} a door. The Earth \emph{pulls} you downwards. A rope \emph{pulls} a sledge. The ground \emph{pushes} up on your feet.
\end{skillbox}

\subsection{What Forces Can Do}\label{sub:1.2}

A force can only change an object in a small number of ways. Learn these five and you can describe almost any situation in this unit.

\vspace{4pt}
\begin{ctab}
\begin{tabular}{|C{0.8cm}|L{5.6cm}|L{8.4cm}|}
  \thead{\#} & \thead{A force can\ldots} & \thead{Example} \\ \hline
  \talt 1 & Make a still object start moving   & A kicked football leaves the ground \\ \hline
        2 & Make a moving object stop          & Brakes bring a bicycle to a halt \\ \hline
  \talt 3 & Make an object speed up or slow down & Pedalling harder makes you accelerate \\ \hline
        4 & Change the direction of motion     & A tennis racquet sends the ball back \\ \hline
  \talt 5 & Change the shape of an object      & Squashing a drink can \\ \hline
\end{tabular}
\end{ctab}

\subsection{Measuring Force: The Newton}\label{sub:1.3}

Force is measured in \textbf{newtons}, written \textbf{N}. The unit is named after Sir Isaac Newton, whose laws of motion you will meet in Lesson~4.

\begin{infobox}{Getting a feel for a newton}
\textbf{1 N} is roughly the force you feel holding a small apple (about 100 g) against gravity.\\
\textbf{10 N} is about the weight of a 1 kg bag of sugar.\\
\textbf{500 N} is roughly the weight of an average adult.
\end{infobox}

We measure force with a \textbf{spring scale} (also called a \emph{newton meter}). Inside is a spring: the harder you pull, the further it stretches, and a pointer moves down a scale marked in newtons.

\vspace{4pt}
\diagSpringScales
\figcap[fig:spring-scale]{Reading a spring scale. Always read at eye level, and check the size of each small division first.}

\begin{skillbox}{How to read a spring scale correctly}
\begin{enumerate}[itemsep=1pt, leftmargin=1.4em, topsep=1pt]
  \item Check the \textbf{largest} number on the scale --- this is the maximum it can measure.
  \item Work out what \textbf{one small division} is worth. If there are 10 small gaps between 0 and 5 N, each is 0.5 N.
  \item Hold the scale still and read at \textbf{eye level} to avoid parallax error.
  \item Read the value where the \textbf{pointer} sits, and always write the \textbf{unit (N)}.
\end{enumerate}
\end{skillbox}

\begin{workedex}{Reading a spring scale}
A spring scale is marked from 0 to 10\,N, with \textbf{20} small divisions between 0 and 10. The pointer sits on the \textbf{7th} small division. What is the reading?

\medskip
\textbf{Step 1 --- one division:}\quad $10 \div 20 = 0.5$\,N per division

\textbf{Step 2 --- the reading:}\quad $7 \times 0.5 = \mathbf{3.5\ \text{N}}$
\end{workedex}

\subsection{Mass and Weight Are Not the Same Thing}\label{sub:1.4}

In everyday speech we use these words as if they mean the same thing. In science they are quite different, and mixing them up causes a lot of lost marks.

\vspace{4pt}
\begin{ctab}
\begin{tabular}{|L{3.2cm}|L{6.0cm}|L{6.0cm}|}
  \thead{} & \thead{Mass} & \thead{Weight} \\ \hline
  \talt What it is   & The amount of matter in an object & The \emph{force} of gravity pulling on that mass \\ \hline
        Unit         & kilograms (kg)                    & newtons (N) \\ \hline
  \talt Measured with& A balance                         & A spring scale \\ \hline
        Changes?     & Same everywhere in the universe   & Changes with the strength of gravity \\ \hline
\end{tabular}
\end{ctab}

\begin{infobox}{The link between them}
\[ \text{Weight} = \text{mass} \times \text{gravitational field strength} \qquad\Rightarrow\qquad \boxed{W = mg} \]
On Earth, $g = 9.8$\,N/kg. On the Moon, $g = 1.6$\,N/kg.
\end{infobox}

\begin{workedex}{Mass and weight}
A student has a mass of 50\,kg. Find their weight on Earth, and on the Moon.

\medskip
\textbf{On Earth:}\quad $W = mg = 50 \times 9.8 = \mathbf{490\ \text{N}}$

\textbf{On the Moon:}\quad $W = mg = 50 \times 1.6 = \mathbf{80\ \text{N}}$

\medskip
Their \emph{mass} is 50\,kg in both places --- only the \emph{weight} changes.
\end{workedex}

\subsection{Contact and Non-contact Forces}\label{sub:1.5}

Forces fall into two groups. \textbf{Contact forces} only act when two objects touch. \textbf{Non-contact forces} act across a gap --- the objects do not need to touch at all.

\vspace{4pt}
\diagForceTypes
\figcap{Contact and non-contact forces}

\begin{workedex}{Classifying forces}
A fridge magnet holds a shopping list on the fridge door without sliding down. Name two forces acting on the magnet and classify each one.

\medskip
\textbf{Weight} pulls the magnet down --- gravity acts across a gap, so it is a \textbf{non-contact} force.

\textbf{Friction} between the magnet and the door stops it sliding --- the surfaces must touch, so it is a \textbf{contact} force.

\medskip
\textit{(The magnetic pull towards the door is also non-contact --- it is what presses the magnet on hard enough for friction to hold it.)}
\end{workedex}

% ─── PART C ───────────────────────────────────────────────────────────────────
\lessonpart{C}{Practice}{Amber}\label{prac:1}

Work through every task. Show your working for calculations, and always include the unit.

\needspace{10\baselineskip}
\task For each everyday action, decide whether the force is a push or a pull. Write \textbf{PUSH} or \textbf{PULL} in the box.

\vspace{4pt}
\begin{ctab}
\begin{tabular}{|SW{5.2cm}|Z{2.3cm}|W{5.2cm}|Z{2.3cm}|}
  \hd & \thead{Action} & \thead{Push or pull?} & \thead{Action} & \thead{Push or pull?} \\ \hline
  & Kicking a football      & \scaf{PUSH} & Opening a drawer       & \scaf{PULL} \\ \hline
  & A crane lifting a girder & & Squeezing a stress ball & \\ \hline
  & A dog straining on a lead & & Wind on a sail       & \\ \hline
\end{tabular}
\end{ctab}

\needspace{10\baselineskip}
\task For each event below, write the \textbf{number} of the effect the force is having. Some events have more than one.

\vspace{2pt}
\begin{ctab}
\begin{tabular}{|C{2.6cm}|C{2.6cm}|C{3.0cm}|C{2.6cm}|C{2.6cm}|}
  \thead{1} & \thead{2} & \thead{3} & \thead{4} & \thead{5} \\ \hline
  starts it moving & stops it & speeds it up or slows it down & changes its direction & changes its shape \\ \hline
\end{tabular}
\end{ctab}

\vspace{3pt}
\begin{ctab}
\begin{tabular}{|SW{5.6cm}|Z{1.9cm}|W{5.6cm}|Z{1.9cm}|}
  \hd & \thead{Event} & \thead{Effect (1--5)} & \thead{Event} & \thead{Effect (1--5)} \\ \hline
  & A goalkeeper catches a ball  & \scaf{2} & Stretching a rubber band & \scaf{5} \\ \hline
  & A car turns a corner at steady speed & & A rocket lifts off & \\ \hline
\end{tabular}
\end{ctab}

\needspace{14\baselineskip}
\task Think of three forces you used this morning. For each one, describe the action, write \textbf{PUSH} or \textbf{PULL}, and give the number of the effect it had: \textbf{1} starts it moving, \textbf{2} stops it, \textbf{3} speeds it up or slows it down, \textbf{4} changes its direction, \textbf{5} changes its shape.

\vspace{3pt}
\begin{ctab}
\begin{tabular}{|SW{8.4cm}|Z{2.8cm}|Z{2.8cm}|}
  \hd & \thead{Action} & \thead{Push or pull?} & \thead{Effect (1--5)} \\ \hline
  & \scaf{Opening the fridge door} & \scaf{PULL} & \scaf{1} \\ \hline
  & & & \\ \hline
  & & & \\ \hline
  & & & \\ \hline
\end{tabular}
\end{ctab}

\needspace{24\baselineskip}
\task Write the reading shown on each of the three spring scales below. Check what one small division is worth on each scale first.

\diagSpringScales

\vspace{3pt}
\begin{ctab}
\begin{tabular}{|SZ{4.6cm}|Z{4.6cm}|Z{4.6cm}|}
  \hd & \thead{Scale A} & \thead{Scale B} & \thead{Scale C} \\ \hline
  & & & \\ \hline
\end{tabular}
\end{ctab}

\needspace{10\baselineskip}
\task Complete the table. Use $g = 9.8$\,N/kg on Earth and $g = 1.6$\,N/kg on the Moon.

\vspace{3pt}
\begin{ctab}
\begin{tabular}{|SW{4.6cm}|Z{2.8cm}|Z{3.6cm}|Z{3.6cm}|}
  \hd & \thead{Object} & \thead{Mass (kg)} & \thead{Weight on Earth (N)} & \thead{Weight on Moon (N)} \\ \hline
  & Bag of sugar & 1   & \scaf{$1 \times 9.8 = 9.8$} & \scaf{$1 \times 1.6 = 1.6$} \\ \hline
  & Bicycle      & 12  & & \\ \hline
  & Small car    & 900 & & \\ \hline
\end{tabular}
\end{ctab}

\taskpair{A bag of flour has a mass of 2\,kg. Find its weight on Earth. Use $g = 9.8$\,N/kg.}{2.2cm}
         {Which has the greater weight on Earth: a 5\,kg dog, or a box that weighs 50\,N? Show your working.}{2.2cm}

\needspace{10\baselineskip}
\task Complete the table. For each force, write whether it is a contact or non-contact force, and give a real-world example.

\vspace{4pt}
\begin{ctab}
\begin{tabular}{|SW{3.4cm}|Z{3.4cm}|W{7.8cm}|}
  \hd & \thead{Force} & \thead{Contact / Non-contact} & \thead{Real-world example} \\ \hline
  & Gravity        & \scaf{Non-contact} & \scaf{An apple falling from a tree} \\ \hline
  & Friction       & & \\ \hline
  & Magnetic force & & \\ \hline
  & Normal force   & & \\ \hline
  & Air resistance & & \\ \hline
  & Tension        & & \\ \hline
\end{tabular}
\end{ctab}


\needspace{10\baselineskip}
\task A spring scale is marked from 0 to 20\,N, with \textbf{40} small divisions. The pointer sits on the \textbf{13th} small division. What is the reading? Show both steps.
\answerbox{2.2cm}

\taskpair{A cat has a mass of 4\,kg. Find its weight on Earth.}{2.2cm}
         {The same cat travels to the Moon. Find its weight there. Has its mass changed?}{2.2cm}

\needspace{10\baselineskip}
\task A balloon rubbed on your hair sticks to a wall. Name the force holding it there, and say whether it is a contact or non-contact force.
\answerbox{2.0cm}

% ─── PART D ───────────────────────────────────────────────────────────────────
\lessonpart{D}{Above level}{Violet}\label{ext:1}

\begin{yr8box}
These questions go beyond the Year 7 standard. They ask you to combine ideas, use a formula, or judge someone else's reasoning.

\medskip
\needspace{10\baselineskip}
\exttask An astronaut has a mass of 75\,kg.
\pt{a}Find her weight on Earth.
\pt{b}Find her weight on the Moon.
\pt{c}Explain why she can jump much higher on the Moon, even though her mass has not changed.
\answerbox{3.4cm}


\needspace{10\baselineskip}
\exttask A suitcase has a weight of 196\,N on Earth.
\pt{a}Find its mass.
\pt{b}Find its weight on the Moon.
\answerbox{3.0cm}

\needspace{10\baselineskip}
\exttask A bathroom scale really measures \emph{force}, but it is labelled in kilograms. A 60\,kg student reads 60\,kg on it on Earth. Predict what the same scale would read on the Moon, and explain why the scale is not really measuring mass.
\answerbox{3.0cm}

\needspace{10\baselineskip}
\exttask On Mars, $g = 3.7$\,N/kg. A 10\,kg backpack is carried on Earth ($g = 9.8$\,N/kg), on the Moon ($g = 1.6$\,N/kg) and on Mars.
\pt{a}Calculate the backpack's weight in each place.
\pt{b}Rank the three places from easiest to hardest to carry the backpack, and explain your ranking.
\answerbox{3.4cm}

\needspace{10\baselineskip}
\exttask A rock hangs from a spring scale on Earth, and the scale reads 4.9\,N. Use $g = 9.8$\,N/kg on Earth and $g = 1.6$\,N/kg on the Moon.
\pt{a}Find the mass of the rock.
\pt{b}Find what the same spring scale would read if the rock were hung from it on the Moon.
\answerbox{3.2cm}

\needspace{10\baselineskip}
\exttask A \emph{balance} compares an object with known masses. A \emph{spring scale} measures the pull of gravity on the object. Explain why a balance gives the same result on the Moon as on Earth, but a spring scale does not.
\answerbox{3.2cm}

\needspace{10\baselineskip}
\exttask Magnets can push or pull without touching. Describe a quick test you could do with a magnet, a paperclip and a sheet of paper to show that magnetic force acts \emph{through} the paper. Say what you would do and what you would expect to see.
\answerbox{3.4cm}

\end{yr8box}
\newpage
% ══════════════════════════════════════════════════════════════════════════════
%  LESSON 2 — FORCE DIAGRAMS AND NET FORCE                        (Week 1 · P2)
% ══════════════════════════════════════════════════════════════════════════════
\sectiondiv{Force Diagrams and Net Force}{Naming forces \textbullet{} Force arrows \textbullet{} Free body diagrams \textbullet{} Balanced and unbalanced forces}\label{sec:2}
\lessoncodes{VC2S8U14 \textbullet{} VC2S8I04}

% ─── PART A ───────────────────────────────────────────────────────────────────
\lessonpart{A}{Reading and comprehension}{Teal}\label{read:2}

\begin{readingbox}{Reading --- Winning a Tug-of-War}
In a tug-of-war, two teams pull on the same rope in opposite directions. While the rope stays still, the two pulls are exactly the same size. They are \emph{balanced}, so the net force on the rope is zero --- the rope does not move, even though both teams are pulling as hard as they can.

The moment one team pulls even a little harder, the forces become \emph{unbalanced}. There is now a net force towards the stronger team, and the rope starts to move their way. The bigger the difference between the two pulls, the faster the rope speeds up.

So how does a team pull harder? Strength helps, but the real secret is under their feet. To pull on the rope, each person pushes backwards against the ground with their shoes. Friction between the shoes and the ground stops them from sliding forwards. A team with good grip --- rough soles, a dry surface, bodies low and leaning back --- can pull harder before their feet slip. A team on a slippery surface will be dragged across, no matter how strong they are.

Coaches explain this with force arrows. They draw the rope, the two pulls and the friction on the feet, each as an arrow whose length shows its size. One look at the arrows shows which team is going to win.
\end{readingbox}

\rcreset
\rcq{While the rope is not moving, what can you say about the two teams' pulls? Use the word \emph{balanced}.}{1.6cm}
\rcq{What happens to the rope when one team pulls harder? Explain why, using the words \emph{net force}.}{1.8cm}
\rcq{According to the passage, why does a team on a slippery surface lose, however strong they are?}{1.8cm}
\rcq{List two things from the passage that give a team better grip.}{1.4cm}
\rcq{How can a coach tell which team will win just by looking at the force arrows?}{1.6cm}
\rcq{Explain why each person needs friction under their shoes in order to pull on the rope.}{1.8cm}
\rcq{A team plays in socks on a polished wooden floor. Use the passage to predict what will happen to them, and explain why.}{1.8cm}
\rcq{In your own words, explain the difference between balanced and unbalanced forces. Give an example of each from the passage.}{2.2cm}

% ─── PART B ───────────────────────────────────────────────────────────────────
\lessonpart{B}{Theory and examples}{Blue}\label{theory:2}

\subsection{Naming the Forces}\label{sub:2.1}

Before you can draw a force diagram, you need the right names. These are the forces you will use again and again in this unit.

\vspace{4pt}
\begin{ctab}
\begin{tabular}{|L{3.4cm}|L{6.4cm}|L{5.4cm}|}
  \thead{Force} & \thead{What causes it} & \thead{Which way it acts} \\ \hline
  \talt Weight ($W$)        & Gravity pulling on mass              & Always straight down \\ \hline
        Normal force ($N$)  & A surface pushing back on an object  & Perpendicular to the surface \\ \hline
  \talt Friction ($f$)      & Two surfaces rubbing together        & Opposite to the motion \\ \hline
        Air resistance      & Air pushing against a moving object  & Opposite to the motion \\ \hline
  \talt Tension ($T$)       & A rope, string or cable pulling      & Along the rope, away from the object \\ \hline
        Applied force       & A person or machine pushing/pulling  & Whichever way it is applied \\ \hline
\end{tabular}
\end{ctab}

\subsection{Drawing Force Arrows --- The Three Rules}\label{sub:2.2}

\begin{skillbox}{The three rules for every force arrow}
\begin{enumerate}[itemsep=2pt, leftmargin=1.4em, topsep=1pt]
  \item \textbf{Start the arrow on the object.} The tail sits on the object the force acts on.
  \item \textbf{Point the arrowhead the way the force acts.} Down for weight, up for the normal force, and so on.
  \item \textbf{Make the length match the size.} A 20\,N force is drawn twice as long as a 10\,N force. \emph{Always}.
\end{enumerate}
\end{skillbox}

Rule 3 is the one students most often forget --- and it matters, because the picture should tell you the answer at a glance. If two arrows are the same length, the forces are balanced. If one is longer, that is the direction the object accelerates.

\vspace{4pt}
\diagArrowScale
\figcap{Using a scale for force arrows. Here 1\,cm represents 5\,N, so a 15\,N force is drawn 3\,cm long.}

\begin{workedex}{Drawing arrows to scale}
Using a scale of 1\,cm = 5\,N, how long should the arrow be for a 25\,N push?

\medskip
\textbf{Length} $=$ force $\div$ scale $= 25 \div 5 = \mathbf{5\ \text{cm}}$
\end{workedex}

\subsection{Free Body Diagrams}\label{sub:2.3}

A \textbf{free body diagram (FBD)} shows \emph{all} the forces acting on \emph{one} object, each drawn as a labelled arrow using the three rules. The object is drawn as a simple box or dot --- nothing else from the picture is included.

\diagFBD
\figcap[fig:fbd]{Balanced forces (left) and unbalanced forces (right)}

\begin{workedex}{Drawing a free body diagram}
A book with a weight of 10\,N rests on a table. Draw its free body diagram.

\medskip
\textbf{Step 1 --- find every force.} Gravity pulls the book down (\emph{weight}). The table pushes up on it (\emph{normal force}).

\textbf{Step 2 --- size each force.} The book is not moving up or down, so the table pushes up exactly as hard as gravity pulls down: both are 10\,N.

\textbf{Step 3 --- draw.} A box for the book, a 10\,N arrow down labelled $W$, and an arrow of the \emph{same length} up labelled $N$ --- as in the left of Figure~\ref{fig:fbd}.
\end{workedex}

\subsection{Balanced and Unbalanced: Net Force}\label{sub:2.4}

When you add up all the forces on an object, the single force left over is called the \textbf{net force} (or resultant force).

\begin{infobox}{Working out the net force}
\textbf{Same direction} $\rightarrow$ \textbf{add} the forces.\\
\textbf{Opposite directions} $\rightarrow$ \textbf{subtract} the smaller from the larger. The net force acts in the direction of the larger force.\\[3pt]
If the net force is \textbf{zero}, the forces are \textbf{balanced} and the motion does not change.\\
If the net force is \textbf{not zero}, the forces are \textbf{unbalanced} and the object speeds up, slows down or changes direction.
\end{infobox}

\begin{workedex}{Finding the net force}
A sledge is pulled forwards with 45\,N. Friction pushes back with 18\,N. Find the net force.

\medskip
Opposite directions, so subtract: \quad $45 - 18 = \mathbf{27\ \text{N forwards}}$

\medskip
The net force is not zero, so the forces are unbalanced and the sledge \textbf{speeds up}.
\end{workedex}

% ─── PART C ───────────────────────────────────────────────────────────────────
\lessonpart{C}{Practice}{Amber}\label{prac:2}

Work through every task. Show your working for calculations, and always include the unit.

\needspace{10\baselineskip}
\task Using a scale of \textbf{1\,cm = 2\,N}, how long should the arrow be for each force?
\pt{a}6\,N
\pt{b}10\,N
\pt{c}3\,N
\answerbox{2.0cm}

\taskpair{Using a scale of 1\,cm = 10\,N, how long should the arrow be for a 45\,N pull?}{2.0cm}
         {An arrow is 3.5\,cm long. The scale is 1\,cm = 20\,N. What force does it show?}{2.0cm}

\needspace{14\baselineskip}
\task A box is pulled right with 30\,N and friction resists with 10\,N. Using a scale of \textbf{1\,cm = 10\,N}, draw both arrows on the box below. Label each with its name and size.

\vspace{3pt}
\begin{tcolorbox}[enhanced, colback=white, colframe=MidGrey, height=4.2cm, boxrule=0.5pt]
\end{tcolorbox}

\needspace{24\baselineskip}
\task[a] The diagrams below show the forces on a book resting on a table (left) and on a box being pushed along the floor (right).

\diagFBD
Why is the book not speeding up or slowing down? Use the words \textit{net force} in your answer.
\answerbox{2.2cm}

\taskpart{b} The box is pushed right with 20\,N, and friction pushes left with 5\,N. Calculate the net force on the box and state which way it speeds up.
\answerbox{2.2cm}

\needspace{16\baselineskip}
\taskpart{c} Draw a free body diagram for a skydiver falling at a steady speed (terminal velocity). Label each force with its name and show their comparative sizes.

\vspace{3pt}
\drawbox{4.6cm}{
\textit{\color{MidGrey}Draw your force diagram here.}
}

\needspace{14\baselineskip}
\task A lamp with a weight of 15\,N hangs still from the ceiling on a cord. Draw its free body diagram, labelling each force with its name and size.

\vspace{3pt}
\drawbox{4.2cm}{
\textit{\color{MidGrey}Draw your force diagram here.}
}

\needspace{6\baselineskip}
\textbf{Net force.} For each question, state the size \emph{and} direction of the net force, and whether the forces are balanced or unbalanced.

\taskpair{Two people push a car in the same direction with 250\,N and 310\,N. Find the net force.}{2.2cm}
         {A trolley is pushed right with 40\,N while friction pushes left with 12\,N. Find the net force.}{2.2cm}

\taskpair{A book rests on a table. Its weight is 8\,N. What is the normal force from the table? Explain.}{2.2cm}
         {In a tug-of-war, one team pulls left with 900\,N and the other pulls right with 900\,N. Describe the motion of the rope.}{2.2cm}

\needspace{10\baselineskip}
\task A rocket engine produces 50\,000\,N of upward thrust. The rocket's weight is 32\,000\,N.
\pt{a}Find the net force.
\pt{b}State which way the rocket speeds up.
\pt{c}Explain what would happen if the thrust dropped to 32\,000\,N.
\answerbox{3.2cm}

\taskpair{A box has three forces on it: 15\,N right, 10\,N right and 30\,N left. Find the size and direction of the net force.}{2.4cm}
         {A skydiver weighs 600\,N. Air resistance on her is 600\,N. Are the forces balanced or unbalanced? Describe her motion.}{2.4cm}

\needspace{10\baselineskip}
\task A swimmer moves forward at a steady speed. Her arms provide 120\,N of forward force.
\pt{a}What is the size of the water resistance acting on her?
\pt{b}Explain how you know.
\answerbox{2.6cm}

\needspace{10\baselineskip}
\task A parachutist has a weight of 700\,N. Just after the parachute opens, air resistance is 950\,N.
\pt{a}Find the net force and its direction.
\pt{b}Describe what happens to the parachutist's speed.
\pt{c}Air resistance then falls to 700\,N --- describe the motion now.
\answerbox{3.4cm}


% ─── PART D ───────────────────────────────────────────────────────────────────
\needspace{12\baselineskip}
\task In a tug-of-war, Team A pulls the rope left with 550\,N and Team B pulls it right with 620\,N.
\pt{a}Find the size and direction of the net force on the rope.
\pt{b}Which team is winning? Explain using the words \emph{unbalanced} and \emph{net force}.
\pt{c}Draw both pulls on the rope below as force arrows, using a scale of 1\,cm = 100\,N. Label each arrow with its size.

\vspace{3pt}
\drawbox{4.4cm}{
\vspace*{2.4cm}
\begin{center}\tikz{\draw[line width=2pt, DarkGrey] (0,0) -- (6,0);}\\[2pt]\textit{\small\color{MidGrey}rope}\end{center}
}

\lessonpart{D}{Above level}{Violet}\label{ext:2}

\begin{yr8box}
These questions go beyond the Year 7 standard. They ask you to combine ideas, use a formula, or judge someone else's reasoning.

\medskip
\needspace{10\baselineskip}
\exttask Sam says: \emph{``A hockey puck sliding across ice at a steady speed must have an unbalanced force pushing it forwards.''} Is Sam right? Use a free body diagram and the idea of net force to explain.
\answerbox{3.4cm}

\needspace{14\baselineskip}
\exttask A skydiver jumps from a plane. Draw a free body diagram for each moment below. Then explain why air resistance grows, and why the skydiver eventually stops speeding up.
\pt{i}just after jumping
\pt{ii}while still speeding up
\pt{iii}at terminal velocity

\vspace{3pt}
\edef\abgrp{\number\numexpr\value{abox}+1\relax}%
\noindent\begin{minipage}[t]{0.31\linewidth}\drawbox{4.2cm}{\textbf{(i)} just after jumping}\end{minipage}\hfill
\begin{minipage}[t]{0.31\linewidth}\drawbox{4.2cm}{\textbf{(ii)} speeding up}\end{minipage}\hfill
\begin{minipage}[t]{0.31\linewidth}\drawbox{4.2cm}{\textbf{(iii)} terminal velocity}\end{minipage}\par\gdef\abgrp{0}
\answerbox{2.6cm}

\needspace{10\baselineskip}
\exttask A car has a driving force of 4500\,N and a total resistive force of 1200\,N.
\pt{a}Calculate the net force.
\pt{b}State which way the car accelerates.
\pt{c}Suggest one way the driver could make the forces balanced.
\answerbox{2.8cm}

\needspace{10\baselineskip}
\exttask A 1200\,kg car has a driving force of 4800\,N. Air resistance and friction total 1500\,N.
\pt{a}Find the net force.
\pt{b}Use $F = ma$ (acceleration = net force $\div$ mass) to find the acceleration.
\pt{c}The driving force rises to 6000\,N --- find the new acceleration.
\answerbox{3.6cm}
\needspace{16\baselineskip}
\exttask A car drives along a flat road at a steady 60\,km/h. Draw a free body diagram for the car showing four forces: weight, the normal force, the driving force and the total resistive force. Show which pairs of forces are balanced.

\vspace{3pt}
\drawbox{4.6cm}{
\textit{\color{MidGrey}Draw your force diagram here.}
}

\needspace{10\baselineskip}
\exttask A 2\,kg ball is falling. Air resistance pushes up on it with 5\,N. Use $W = mg$ with $g = 9.8$\,N/kg.
\pt{a}Find the weight of the ball.
\pt{b}Find the size and direction of the net force.
\pt{c}Is the ball speeding up, slowing down or moving at a steady speed? Explain.
\answerbox{3.4cm}

\needspace{10\baselineskip}
\exttask A loaded truck and a small car are both travelling at the same speed. Explain, using the idea of net force, why the truck needs a much bigger braking force to stop in the same distance.
\answerbox{3.0cm}

\end{yr8box}
\newpage
% ══════════════════════════════════════════════════════════════════════════════
%  LESSON 3 — EXPERIMENT 1: INVESTIGATING FRICTION                (Week 1 · P3)
% ══════════════════════════════════════════════════════════════════════════════
\sectiondiv{Experiment 1: Investigating Friction}{Prediction \textbullet{} Fair testing \textbullet{} Repeat trials \textbullet{} Means}\label{sec:3}\label{exp:1}
\lessoncodes{VC2S8U14 \textbullet{} VC2S8I01 \textbullet{} VC2S8I02 \textbullet{} VC2S8I03 \textbullet{} VC2S8I04 \textbullet{} VC2S8I05}

% ─── PART A ───────────────────────────────────────────────────────────────────
\lessonpart{A}{Reading and comprehension}{Teal}\label{read:3}

\begin{readingbox}{Reading --- Grip}
Friction is the force that makes walking possible. With every step, your shoe pushes backwards against the ground, and friction pushes you forwards. On ice, where friction is tiny, your foot slips backwards and you go nowhere.

Friction comes from the roughness of surfaces. Even surfaces that look smooth are covered in tiny bumps when you look at them under a microscope. When two surfaces press together, the bumps catch on each other and resist sliding. Rougher surfaces have bigger bumps that catch more, so they produce more friction. Pressing the surfaces together harder also increases friction.

Sometimes engineers want as much friction as possible --- in the rubber of a car tyre, the soles of running shoes or the grip tape on a skateboard. Other times they want as little as possible: oil in an engine and wax on skis reduce friction, so parts slide easily and less energy is wasted as heat.

In this experiment you will measure friction with a spring scale. When you pull a block along at a \emph{steady} speed, your pull and the friction are balanced --- so the reading on the spring scale is equal to the friction force.
\end{readingbox}

\rcreset
\rcq{Explain how friction helps you to walk, and why walking on ice is so hard.}{1.8cm}
\rcq{Why do rougher surfaces produce more friction? Use the word \emph{bumps} in your answer.}{1.6cm}
\rcq{Give one example where engineers want more friction, and one where they want less.}{1.6cm}
\rcq{Why must the block move at a \emph{steady} speed when you read the spring scale?}{1.8cm}

% ─── PART B ───────────────────────────────────────────────────────────────────
\lessonpart{B}{Planning the experiment}{Blue}\label{theory:3}

\begin{expheader}{1}{Investigating Friction}{To investigate how the type of surface affects the friction force on a sliding object.}
\medskip\small\textbf{Curriculum:} VC2S8U14 \textbullet{} VC2S8I01 \textbullet{} VC2S8I02 \textbullet{} VC2S8I03 \textbullet{} VC2S8I04 \textbullet{} VC2S8I05
\end{expheader}

\subsection{Making it a Fair Test}\label{sub:3.1}

In a fair test you change \textbf{one} thing, measure \textbf{one} thing, and keep everything else the same.

\vspace{4pt}
\begin{ctab}
\begin{tabular}{|L{3.6cm}|L{11.4cm}|}
  \thead{Variable} & \thead{What it means} \\ \hline
  \talt Independent & The one thing \emph{you} change on purpose \\ \hline
        Dependent   & The thing you \emph{measure} to see the effect \\ \hline
  \talt Controlled  & Everything you keep the \emph{same}, so the test is fair \\ \hline
\end{tabular}
\end{ctab}

\begin{workedex}{Identifying variables}
A student rolls a toy car down a ramp at different heights and measures how far it travels across the floor. Identify the variables.

\medskip
\textbf{Independent:} the height of the ramp (the student changes it).

\textbf{Dependent:} the distance the car travels (the student measures it).

\textbf{Controlled:} the same car, the same floor surface, releasing the car without pushing it.
\end{workedex}

\begin{scaffold}{Prediction --- before you begin}
Which surface do you think will produce the most friction? Why?
\answerbox{2cm}
\end{scaffold}

\subsubsection*{Equipment}
Wooden block (or heavy book) \quad Spring scale (newton meter) \quad 3+ surfaces (bench, sandpaper, fabric, carpet) \quad Ruler

\subsubsection*{Method}
\begin{enumerate}[itemsep=2pt]
  \item Place the wooden block on the first surface.
  \item Attach the spring scale to the block. Pull slowly and steadily.
  \item Read the spring scale while the block moves at a steady speed. Record the friction force in newtons.
  \item Repeat 3 times for this surface and record all values.
  \item Repeat steps 1--4 for each surface.
\end{enumerate}

\begin{infobox}{Safety}
Pull gently and keep the scale in line with the block --- a spring scale that is overstretched can snap back. Keep sandpaper away from eyes and skin.
\end{infobox}

% ─── PART C ───────────────────────────────────────────────────────────────────
\lessonpart{C}{Results and analysis}{Amber}\label{prac:3}

\subsubsection*{Results}

\begin{ctab}
\begin{tabular}{|SW{3.6cm}|Z{2.4cm}|Z{2.4cm}|Z{2.4cm}|Z{2.8cm}|}
  \hd & \thead{Surface} & \thead{Trial 1 (N)} & \thead{Trial 2 (N)} & \thead{Trial 3 (N)} & \thead{Mean (N)} \\ \hline
  & \scaf{e.g.\ Bench} & \scaf{1.2} & \scaf{1.4} & \scaf{1.3} & \scaf{1.3} \\ \hline
  & & & & & \\ \hline
  & & & & & \\ \hline
  & & & & & \\ \hline
  & & & & & \\ \hline
\end{tabular}
\end{ctab}

\begin{infobox}{Calculating a mean}
Add the trials, then divide by how many there are. For 1.2\,N, 1.4\,N and 1.3\,N: \quad $(1.2 + 1.4 + 1.3) \div 3 = 3.9 \div 3 = \mathbf{1.3\ \text{N}}$
\end{infobox}

\needspace{10\baselineskip}
\task Identify the independent variable, the dependent variable, and two controlled variables in \emph{your} experiment.
\answerbox{2.6cm}

\needspace{10\baselineskip}
\task Which surface produced the greatest friction? Which produced the least? Was your prediction correct?
\answerbox{2.2cm}

\needspace{10\baselineskip}
\task Why did you take three trials for each surface? What did you do with the three values?
\answerbox{2.2cm}

\taskpair{Another group measured 2.1\,N, 2.3\,N and 2.2\,N on carpet. Calculate the mean.}{2.2cm}
         {A group measured 1.8\,N, 1.9\,N and 2.9\,N on fabric. Which value looks wrong? Find the mean of the other two.}{2.2cm}

\needspace{10\baselineskip}
\task Rough surfaces are covered in bigger bumps, which catch on each other and resist sliding. Use this idea to explain \emph{why} the surface that gave the most friction in your experiment did so.
\answerbox{2.2cm}

% ─── PART D ───────────────────────────────────────────────────────────────────
\needspace{20\baselineskip}
\task Draw a column graph of the \emph{mean} friction force for each surface you tested. Put the surfaces on the horizontal axis and \textbf{friction force (N)} on the vertical axis. Give your graph a title.

\vspace{3pt}
\begin{tcolorbox}[enhanced, colback=white, colframe=MidGrey, boxrule=0.5pt, height=7.5cm,
  underlay={\begin{tcbclipinterior}\draw[step=5mm, Line, very thin] (interior.south west) grid (interior.north east);\end{tcbclipinterior}}]
\end{tcolorbox}

\lessonpart{D}{Above level}{Violet}\label{ext:3}

\begin{yr8box}
These questions go beyond the Year 7 standard. They ask you to evaluate the method, predict, or use a formula.

\medskip
\needspace{10\baselineskip}
\exttask Even on the same surface, your three readings were probably not identical. Suggest one reason why, and one change to the method that would make the readings more consistent.
\answerbox{2.6cm}

\needspace{10\baselineskip}
\exttask Predict what would happen to the friction force if you stacked a second block on top of the first. Explain your prediction.\par\noindent\textit{Hint: pressing two surfaces together harder makes their bumps catch more.}
\answerbox{2.4cm}

\needspace{10\baselineskip}
\exttask The wooden block has a mass of 0.3\,kg. Choose your highest mean friction force. If the block were sliding and you let go of the spring scale, friction would be the only horizontal force on it. Use $F = ma$ (acceleration = force $\div$ mass) to calculate how quickly it would slow down. Show all working.
\answerbox{3.2cm}
\needspace{12\baselineskip}
\exttask Design a fair test to find out whether the \emph{mass} of the block changes the friction force on the bench.
\pt{a}Name the independent and dependent variables.
\pt{b}Name two variables you would keep the same.
\pt{c}Predict what you would find, and give a reason.
\answerbox{4.2cm}

\end{yr8box}

\newpage
% ══════════════════════════════════════════════════════════════════════════════
%  LESSON 4 — NEWTON'S FIRST AND THIRD LAWS                       (Week 2 · P1)
% ══════════════════════════════════════════════════════════════════════════════
\sectiondiv{Newton's First and Third Laws}{Inertia \textbullet{} Why moving things keep moving \textbullet{} Action and reaction \textbullet{} Force pairs}\label{sec:4}
\lessoncodes{VC2S8U14 \textbullet{} VC2S10U17 (above level)}

% ─── PART A ───────────────────────────────────────────────────────────────────
\lessonpart{A}{Reading and comprehension}{Teal}\label{read:4}

\begin{readingbox}{Reading --- Why Rockets Work in Empty Space}
When rockets were first proposed as a way of reaching space, some newspapers scoffed at the idea. In 1920 an editorial in the \emph{New York Times} argued that a rocket could never work beyond the atmosphere, because out in the vacuum there would be nothing for the exhaust gases to push against. It seemed like common sense. A swimmer pushes against water, a runner pushes against the ground --- so surely a rocket must push against the air?

The reasoning was wrong, and the mistake came from misunderstanding Newton's Third Law. A rocket does not push against the air behind it. It pushes against \emph{its own exhaust}. The engine hurls a huge mass of hot gas backwards, and by the Third Law the gas pushes the rocket forwards with exactly the same force. Nothing outside the rocket is needed at all. In fact rockets work \emph{better} in space, because there is no air resistance slowing them down.

Out in deep space, something even more surprising happens: a spacecraft can switch its engines off and simply keep going. Far from any planet there is no air and almost nothing pushing or pulling on it, so there is almost no unbalanced force to slow it down. The two \emph{Voyager} space probes were launched in 1977. No engine has pushed them forwards for more than thirty years, yet they are still travelling away from the Sun at more than 55\,000 kilometres per hour. Here on Earth, a ball rolled along the ground soon stops --- but only because friction and air resistance act on it. Take those away, and a moving object just keeps moving.

The \emph{New York Times} eventually printed a correction --- in July 1969, the day after Apollo 11 launched for the Moon, and forty-nine years after the original editorial. It noted, briefly, that further investigation had confirmed the work of Newton in the seventeenth century, and regretted the error.
\end{readingbox}

\rcreset
\rcq{What mistaken reason did the 1920 editorial give for believing rockets could not work in space?}{1.6cm}
\rcq{Explain, using Newton's Third Law, what a rocket actually pushes against.}{2.0cm}
\rcq{The passage says rockets work \emph{better} in space. Give the reason, and name the force that is absent.}{1.6cm}
\rcq{The \emph{Voyager} probes are still moving fast, even though no engine is pushing them. Explain why they have not slowed down.}{2.0cm}
\rcq{According to the passage, why does a ball rolled along the ground on Earth soon stop?}{1.6cm}
\rcq{Which phrase in the second paragraph tells you how the push of the gas on the rocket compares with the push of the rocket on the gas?}{1.4cm}
\rcq{The editorial's reasoning ``seemed like common sense''. Using the two examples in the passage, explain why it seemed right to many people.}{1.8cm}
\rcq{Imagine a \emph{Voyager} probe drifted into a thin cloud of gas and dust in space. Using the passage, predict what would happen to its speed, and explain why.}{2.0cm}

% ─── PART B ───────────────────────────────────────────────────────────────────
\lessonpart{B}{Theory and examples}{Blue}\label{theory:4}

\subsection{Newton's First Law --- The Law of Inertia}\label{sub:4.1}

In Lesson~2 you learned that when the forces on an object are balanced, its motion does not change. Nearly 340 years ago, Sir Isaac Newton made this the first of his three laws of motion.

\diagNewtonOne
\figcap[fig:newton-one]{Newton's First Law}

\begin{infobox}{First Law}
\textit{``An object at rest stays at rest, and an object in motion stays in motion at constant velocity, unless acted upon by an unbalanced (net) force.''}

\medskip
\textbf{Inertia} is the tendency of an object to resist changes to its motion. Greater mass $=$ greater inertia.
\end{infobox}

The First Law has two halves, and both say the same thing: \emph{without an unbalanced force, nothing about the motion changes.}

\vspace{4pt}
\begin{ctab}
\begin{tabular}{|L{3.6cm}|L{5.2cm}|L{6.2cm}|}
  \thead{If the object is\ldots} & \thead{With no unbalanced force, it\ldots} & \thead{Example} \\ \hline
  \talt At rest & stays at rest & A book on a shelf never starts sliding by itself \\ \hline
        Moving  & keeps moving at the same speed, in the same direction & A space probe coasting with its engines off \\ \hline
\end{tabular}
\end{ctab}

\begin{skillbox}{Think about it --- why do things on Earth seem to stop by themselves?}
For almost 2000 years, people followed the Greek thinker Aristotle, who taught that a moving object needs a force to \emph{keep} it moving. It seems right: a ball rolled along the ground soon stops. Newton showed that the ball stops because \textbf{friction} and \textbf{air resistance} are acting on it. They are unbalanced forces, and they slow it down. With no friction and no air resistance, the ball would roll on forever.
\end{skillbox}

\subsubsection*{Mass and inertia}
The more mass an object has, the more inertia it has --- so the harder it is to change its motion in \emph{any} way.

\vspace{4pt}
\begin{ctab}
\begin{tabular}{|L{4.4cm}|L{5.0cm}|L{5.0cm}|}
  \thead{} & \thead{Empty shopping trolley} & \thead{Full shopping trolley} \\ \hline
  \talt Mass and inertia      & Small & Large \\ \hline
        Getting it moving     & Easy  & Hard \\ \hline
  \talt Stopping it           & Easy  & Hard \\ \hline
        Steering it round a corner & Easy  & Hard \\ \hline
\end{tabular}
\end{ctab}

\begin{workedex}{At rest stays at rest}
A coin sits on a playing card on top of a glass. The card is flicked quickly sideways. Explain why the coin drops into the glass instead of flying off with the card.

\medskip
\textbf{Step 1 --- before the flick.} The coin is at rest. Its weight pulls down and the card pushes up: the forces are balanced.

\textbf{Step 2 --- the flick.} The big force acts on the \emph{card}, not the coin. Only a tiny friction force acts on the coin, for a very short time.

\textbf{Step 3 --- inertia.} The coin's inertia keeps it at rest where it is. With the card gone, nothing pushes up on the coin, so its weight is now \textbf{unbalanced} and it falls straight into the glass.
\end{workedex}

\begin{workedex}{Moving keeps moving}
A car travelling at 50\,km/h crashes into a wall and stops suddenly. Explain what happens to a passenger who is not wearing a seatbelt, and how a seatbelt helps.

\medskip
\textbf{Step 1 --- what the crash does.} The wall exerts a huge unbalanced force on the \emph{car}, so the car stops.

\textbf{Step 2 --- inertia.} No force has acted on the passenger yet, so by the First Law the passenger \textbf{keeps moving forwards at 50\,km/h} --- until they hit the dashboard or windscreen.

\textbf{Step 3 --- the seatbelt.} A seatbelt supplies the unbalanced force that slows the passenger down \emph{with} the car, spread across the strong bones of the chest and hips.
\end{workedex}

\subsection{Newton's Third Law --- Action and Reaction}\label{sub:4.2}

Forces always come in pairs. You cannot push on something without it pushing back on you. Press your hand flat against a wall and you can feel the wall pushing back.

\begin{infobox}{Third Law}
\textit{``For every action force, there is an equal and opposite reaction force.''}

\medskip
Action--reaction pairs \textbf{always act on different objects} --- they do \textbf{not} cancel each other.
\end{infobox}

\diagThirdLaw
\figcap[fig:third-law]{Two Third Law pairs. In each pair the forces are the same size, point in opposite directions and act on different objects.}

\begin{skillbox}{How to find the reaction force}
Write the action as ``\emph{A pushes B}'', then swap the names: the reaction is ``\emph{B pushes A}''. The reaction force is always:
\begin{enumerate}[itemsep=1pt, leftmargin=1.4em, topsep=2pt]
  \item the \textbf{same size} as the action force;
  \item in the \textbf{opposite direction};
  \item acting on the \textbf{other object};
  \item the \textbf{same type} of force --- a push answers a push, and a gravity pull answers a gravity pull.
\end{enumerate}
\end{skillbox}

\vspace{4pt}
\begin{ctab}
\begin{tabular}{|L{7.4cm}|L{7.4cm}|}
  \thead{Action force} & \thead{Reaction force} \\ \hline
  \talt Your foot pushes the ground backwards     & The ground pushes your foot forwards \\ \hline
        A swimmer's hands push the water backwards & The water pushes the swimmer's hands forwards \\ \hline
  \talt A rocket pushes its exhaust gas backwards  & The gas pushes the rocket forwards \\ \hline
        The Earth pulls you down (your weight)     & You pull the Earth up \\ \hline
\end{tabular}
\end{ctab}

\subsubsection*{Why don't the two forces cancel?}
Forces only cancel when they act on the \emph{same} object. A Third Law pair always acts on \emph{two different} objects, so the forces can never cancel. When the skaters in Figure~\ref{fig:third-law} push apart, each skater has one unbalanced force acting on them --- so \emph{both} skaters move, in opposite directions. If one skater has less mass, they move away faster. Lesson~5 shows why.

\begin{workedex}{Describing a reaction force}
A student on roller skates pushes against a wall with a force of 40\,N. Describe the reaction force: name it, and give its size and direction. What happens to the student?

\medskip
\textbf{Action:} the student pushes on the wall, 40\,N, towards the wall.

\textbf{Reaction:} the \textbf{wall pushes on the student}, with \textbf{40\,N}, \textbf{away from the wall}.

\textbf{Motion:} this force acts on the student, and the skate wheels give very little friction, so it is \textbf{unbalanced}. The student rolls backwards, away from the wall.
\end{workedex}

% ─── PART C ───────────────────────────────────────────────────────────────────
\lessonpart{C}{Practice}{Amber}\label{prac:4}

Work through every task. Use the words \emph{inertia}, \emph{unbalanced force}, \emph{action} and \emph{reaction} wherever they fit.

\needspace{12\baselineskip}
\task For each situation, decide which law explains it. Write \textbf{FIRST} or \textbf{THIRD} in the box.

\vspace{4pt}
\begin{ctab}
\begin{tabular}{|SW{5.4cm}|Z{2.0cm}|W{5.4cm}|Z{2.0cm}|}
  \hd & \thead{Situation} & \thead{Law} & \thead{Situation} & \thead{Law} \\ \hline
  & Passengers lurch forwards when a bus brakes & \scaf{FIRST} & A swimmer pushes water back and moves forwards & \scaf{THIRD} \\ \hline
  & A puck keeps sliding across smooth ice & & A balloon flies off when its air rushes out & \\ \hline
  & A tablecloth is whipped away and the plates stay put & & A fire hose pushes back hard on the firefighters holding it & \\ \hline
  & When a car turns sharply right, you slide to the left of your seat & & Pushing off the side of a pool sends you away from the wall & \\ \hline
\end{tabular}
\end{ctab}

\needspace{10\baselineskip}
\task Explain why passengers lurch forwards when a bus suddenly brakes. Use the words \textit{inertia} and \textit{unbalanced force} in your answer.
\answerbox{2.6cm}

\needspace{10\baselineskip}
\task A hockey puck slides across an almost frictionless ice surface.
\pt{a}Describe how the puck moves while it slides.
\pt{b}What will eventually stop it, and why?
\answerbox{2.8cm}

\taskpair{Which has more inertia: an empty shopping trolley or a full one? Explain.}{2.4cm}
         {A truck and a small car travel at the same speed. Why does the truck need much stronger brakes?}{2.4cm}

\taskpair{The same push sends a ball rolling across carpet and across a smooth wooden floor. On which surface does it roll further? Explain.}{2.4cm}
         {When you shake a wet umbrella and then stop it suddenly, the water flies off. Use the idea of inertia to explain why.}{2.4cm}

\needspace{10\baselineskip}
\task When a car is hit hard from behind, it is suddenly pushed forwards. Use the First Law to explain why a passenger's head is jolted backwards, and how a headrest helps.
\answerbox{2.8cm}

\needspace{14\baselineskip}
\task Complete the table. For each action force, describe the reaction force.

\vspace{4pt}
\begin{ctab}
\begin{tabular}{|SW{7.2cm}|W{7.6cm}|}
  \hd & \thead{Action force} & \thead{Reaction force} \\ \hline
  & A boat's paddle pushes the water backwards & \scaf{The water pushes the paddle (and the boat) forwards} \\ \hline
  & A hammer pushes down on a nail & \\ \hline
  & A tennis racquet pushes the ball forwards & \\ \hline
  & A bird's wings push the air downwards & \\ \hline
  & The Earth pulls an apple downwards & \\ \hline
\end{tabular}
\end{ctab}

\needspace{10\baselineskip}
\task A swimmer pushes backward against the water with a force of 80\,N. Describe the reaction force: name it, state its size, and state its direction.
\answerbox{2.4cm}

\needspace{16\baselineskip}
\task Two skaters stand still on ice, facing each other. Skater A pushes skater B with a force of 30\,N. Using a scale of \textbf{1\,cm = 10\,N}, draw both forces of the Third Law pair. Draw each skater as a box, start each arrow on the skater it acts on, and label each arrow with its size.

\vspace{3pt}
\drawbox{4.4cm}{
\textit{\color{MidGrey}Draw your force diagram here.}
}

\needspace{10\baselineskip}
\task When you blow up a balloon and let it go, it zooms around the room.
\pt{a}What does the balloon push on?
\pt{b}Use the Third Law to explain why the balloon moves forwards.
\answerbox{2.8cm}

\needspace{10\baselineskip}
\task A skateboarder jumps forwards off the back of a skateboard, and the skateboard shoots off backwards. Explain why, naming the action force and the reaction force.
\answerbox{2.6cm}

\needspace{10\baselineskip}
\task A student steps forwards off a small rowing boat onto a jetty, and the boat drifts backwards, away from the jetty.
\pt{a}Name the action force and the reaction force.
\pt{b}Explain why the boat moves, even though nobody pushed it with their hands.
\answerbox{2.8cm}

% ─── PART D ───────────────────────────────────────────────────────────────────
\lessonpart{D}{Above level}{Violet}\label{ext:4}

\begin{yr8box}
These questions go beyond the Year 7 standard. They ask you to combine ideas, use both laws together, or judge someone else's reasoning.

\medskip
\needspace{10\baselineskip}
\exttask Jess says: \emph{``The weight of a book on a table and the normal force from the table are a Third Law pair, because they are equal and opposite.''} Is Jess right? Explain, and name the true reaction force to each of the two forces.
\answerbox{2.8cm}

\needspace{12\baselineskip}
\exttask For each pair of forces, decide whether it is a Third Law pair. Write \textbf{YES} or \textbf{NO}, and explain how you decided.
\pt{i}The Earth pulls a book down; the book pulls the Earth up.
\pt{ii}A rope pulls a sledge forwards; friction from the snow pulls the sledge backwards.
\pt{iii}A swimmer pushes the water backwards; the water pushes the swimmer forwards.
\answerbox{3.0cm}

\needspace{10\baselineskip}
\exttask A horse pulls a cart. By the Third Law, the cart pulls back on the horse with an equal force. Liam says: \emph{``So the forces cancel out, and the horse can never get the cart moving.''} Explain Liam's mistake. \textit{(Hint: draw a free body diagram for the horse alone. What pushes the horse forwards?)}
\answerbox{3.0cm}

\needspace{10\baselineskip}
\exttask An astronaut on a spacewalk is floating 10\,m from the Space Station, not moving. Her safety line has come loose, but she is holding a heavy toolbox.
\pt{a}Explain how she can use the Third Law to start moving back towards the station.
\pt{b}Once she is moving, does she need to keep doing anything to keep moving? Use the First Law.
\answerbox{3.0cm}

\needspace{10\baselineskip}
\exttask An apple falls from a tree. The Earth pulls the apple down with a force of 1\,N.
\pt{a}What is the reaction force? Give its size and direction, and say what it acts on.
\pt{b}The two forces are the same size. So why do we see the apple move towards the Earth, and not the Earth move towards the apple? Use the idea of inertia.
\answerbox{2.8cm}

\needspace{10\baselineskip}
\exttask Two skaters stand still on ice and push each other apart. Skater A has twice the mass of skater B.
\pt{a}Compare the sizes of the forces on the two skaters.
\pt{b}Which skater moves away faster? Explain, using the idea of inertia.
\answerbox{2.8cm}

\needspace{10\baselineskip}
\exttask A waiter carries a full cup of water on a tray.
\pt{a}The waiter suddenly moves the tray forwards. Which way does the water spill? Explain.
\pt{b}Later, walking at a steady speed, the waiter suddenly stops. Which way does the water spill now? Explain.
\pt{c}Use the First Law to explain why moving the tray smoothly stops the water spilling.
\answerbox{3.4cm}
\end{yr8box}

\newpage
% ══════════════════════════════════════════════════════════════════════════════
%  LESSON 5 — NEWTON'S SECOND LAW                                 (Week 2 · P2)
% ══════════════════════════════════════════════════════════════════════════════
\sectiondiv{Newton's Second Law: $F = ma$}{Acceleration \textbullet{} Force, mass and acceleration \textbullet{} Rearranging $F = ma$ \textbullet{} Using the net force}\label{sec:5}
\lessoncodes{VC2S8U14 \textbullet{} VC2S10U17 (above level)}

% ─── PART A ───────────────────────────────────────────────────────────────────
\lessonpart{A}{Reading and comprehension}{Teal}\label{read:5}

\begin{readingbox}{Reading --- Pushing a Broken-Down Car}
A car has broken down on a quiet, flat street, and it needs to be pushed to the side of the road. One person leans against the back and pushes as hard as they can. Slowly --- very slowly --- the car starts to roll. Two friends run over to help, and with three people pushing, the car speeds up much more quickly. A bigger force gives a bigger acceleration.

Further down the street, a removalist's van packed full of furniture has also broken down. The same three people push it just as hard as they pushed the car. This time it barely creeps forwards. The force is the same, but the van has far more mass, and more mass means a smaller acceleration.

Once the car is rolling, something changes: the pushers can ease off. They only need to push hard enough to match the friction acting on the car, and it keeps rolling at a steady speed. Pushing harder only matters if they want it to go faster.

Isaac Newton tied all of this together in his Second Law: how quickly an object speeds up depends on the net force acting on it, and on its mass. That is why racing cars and racing bikes are built to be as light as possible. With the same engine, or the same legs, a lighter vehicle speeds up faster.
\end{readingbox}

\rcreset
\rcq{What happened to the car's acceleration when two friends joined in? Explain why.}{1.6cm}
\rcq{The same push was used on the loaded van. Describe what happened, and explain it using the word \emph{mass}.}{1.8cm}
\rcq{Once the car rolls at a steady speed, which force are the pushers balancing? How do you know the forces are balanced?}{1.8cm}
\rcq{Why are racing cars and racing bikes built to be as light as possible?}{1.6cm}
\rcq{According to the passage, what \emph{two} things decide how quickly an object speeds up?}{1.4cm}
\rcq{Which sentence in the first paragraph sums up what happened when more people pushed the car?}{1.4cm}
\rcq{Describe one other everyday example where the same push gives a smaller acceleration because the object has more mass.}{1.6cm}
\rcq{Suppose the three pushers kept pushing just as hard after the car was already rolling. Using the passage, predict what would happen to the car, and explain why using the words \emph{net force}.}{2.0cm}

% ─── PART B ───────────────────────────────────────────────────────────────────
\lessonpart{B}{Theory and examples}{Blue}\label{theory:5}

\subsection{What is Acceleration?}\label{sub:5.1}

\textbf{Acceleration} is how quickly an object's speed changes. It is measured in \textbf{metres per second squared (m/s$^2$)}, which means ``metres per second, every second''.

\begin{infobox}{Reading an acceleration}
An acceleration of \textbf{3\,m/s$^2$} means the speed goes up by \textbf{3\,m/s every second}. Starting from rest:

\medskip
\begin{tabular}{|l|c|c|c|c|}
  \hline
  \textbf{Time (s)}    & 0 & 1 & 2 & 3 \\ \hline
  \textbf{Speed (m/s)} & 0 & 3 & 6 & 9 \\ \hline
\end{tabular}

\medskip
Slowing down is also an acceleration. It is sometimes called a \textbf{deceleration}.
\end{infobox}

\subsection{Force, Mass and Acceleration}\label{sub:5.2}

Newton's Second Law links three things: the \textbf{net force} on an object, its \textbf{mass}, and its \textbf{acceleration}.

\diagFma
\figcap[fig:fma]{Newton's Second Law: force, mass and acceleration}

\vspace{2pt}
\begin{ctab}
\begin{tabular}{|L{5.4cm}|L{3.2cm}|L{6.4cm}|}
  \thead{If you\ldots} & \thead{the acceleration\ldots} & \thead{Example} \\ \hline
  \talt increase the force (same mass) & increases & Three people push the car instead of one \\ \hline
        increase the mass (same force) & decreases & The same push on a loaded van \\ \hline
  \talt double the force (same mass)   & doubles   & Two people push a trolley, each as hard as one did before \\ \hline
        double the mass (same force)   & halves    & The same push on two identical trolleys joined together \\ \hline
\end{tabular}
\end{ctab}

\needspace{11\baselineskip}
\begin{infobox}{Second Law}
The acceleration of an object is directly proportional to the net force acting on it, and inversely proportional to its mass.
\[
  \boxed{F = m \times a}
  \qquad\qquad
  \boxed{a = \frac{F}{m}}
  \qquad\qquad
  \boxed{m = \frac{F}{a}}
\]
Force in newtons (N) \quad Mass in kilograms (kg) \quad Acceleration in m/s$^2$
\end{infobox}

\subsection{Using $F = ma$}\label{sub:5.3}

You only need to remember one formula. The \textbf{formula triangle} shows how to rearrange it.

\diagFmaTriangle
\figcap{The formula triangle. Cover the quantity you want; what is left shows how to find it.}

\begin{skillbox}{Four steps for every $F = ma$ question}
\begin{enumerate}[itemsep=1pt, leftmargin=1.4em, topsep=1pt]
  \item Write down what you \textbf{know}, with units.
  \item Choose the \textbf{version} of the formula you need (use the triangle).
  \item \textbf{Substitute} the numbers.
  \item \textbf{Calculate}, and always write the \textbf{unit}.
\end{enumerate}
\end{skillbox}

\begin{workedex}{Finding acceleration}
A 5\,kg trolley is pushed with a net force of 20\,N. What is its acceleration?

\medskip
\textbf{Know:}\quad $F = 20$\,N, \quad $m = 5$\,kg

\textbf{Formula:}\quad $F = m \times a$ \quad $\Rightarrow$ \quad $a = F \div m$

\textbf{Substitute:}\quad $a = 20 \div 5 = \mathbf{4\ \text{m/s}^2}$
\end{workedex}

\begin{workedex}{Finding force}
A 1200\,kg car accelerates at 2\,m/s$^2$. What net force is acting on it?

\medskip
\textbf{Know:}\quad $m = 1200$\,kg, \quad $a = 2$\,m/s$^2$

\textbf{Formula:}\quad $F = m \times a$

\textbf{Substitute:}\quad $F = 1200 \times 2 = \mathbf{2400\ \text{N}}$
\end{workedex}

\begin{workedex}{Finding mass}
A net force of 60\,N gives a box an acceleration of 4\,m/s$^2$. What is the mass of the box?

\medskip
\textbf{Know:}\quad $F = 60$\,N, \quad $a = 4$\,m/s$^2$

\textbf{Formula:}\quad $F = m \times a$ \quad $\Rightarrow$ \quad $m = F \div a$

\textbf{Substitute:}\quad $m = 60 \div 4 = \mathbf{15\ \text{kg}}$
\end{workedex}

\subsection{Always Use the Net Force}\label{sub:5.4}

The $F$ in $F = ma$ is the \textbf{net force} --- the force left over after you add up all the forces (Lesson~2). If the forces are balanced, the net force is zero, so the acceleration is zero: the object stays still, or keeps moving at a steady speed. That is exactly what the First Law says.

\begin{workedex}{Net force first}
A 10\,kg trolley is pushed forwards with 50\,N. Friction pushes back with 10\,N. Find its acceleration.

\medskip
\textbf{Step 1 --- net force:}\quad opposite directions, so subtract: \quad $50 - 10 = 40$\,N forwards

\textbf{Step 2 --- acceleration:}\quad $a = F \div m = 40 \div 10 = \mathbf{4\ \text{m/s}^2}$ \textbf{forwards}
\end{workedex}

% ─── PART C ───────────────────────────────────────────────────────────────────
\lessonpart{C}{Practice}{Amber}\label{prac:5}

Work through every task. Show your working for calculations, and always include the unit.

\needspace{12\baselineskip}
\task For each change, describe what happens to the acceleration. Write \textbf{increases}, \textbf{decreases}, \textbf{doubles} or \textbf{halves}.

\vspace{4pt}
\begin{ctab}
\begin{tabular}{|SW{5.2cm}|Z{2.3cm}|W{5.2cm}|Z{2.3cm}|}
  \hd & \thead{Change} & \thead{Acceleration} & \thead{Change} & \thead{Acceleration} \\ \hline
  & The force is doubled; the mass stays the same & \scaf{doubles} & The mass is doubled; the force stays the same & \scaf{halves} \\ \hline
  & More people help push the same car & & A heavier rider pedals the same bike just as hard & \\ \hline
  & A truck is loaded with bricks; its engine force stays the same & & A rocket burns fuel and gets lighter; its thrust stays the same & \\ \hline
\end{tabular}
\end{ctab}

\needspace{12\baselineskip}
\task Use $F = ma$ to complete the table. The first row is done for you.

\vspace{4pt}
\begin{ctab}
\begin{tabular}{|SW{4.4cm}|Z{3.2cm}|Z{3.2cm}|Z{3.6cm}|}
  \hd & \thead{Situation} & \thead{Force (N)} & \thead{Mass (kg)} & \thead{Acceleration (m/s$^2$)} \\ \hline
  & Trolley pushed      & 20     & 5    & \scaf{$a = 20 \div 5 = 4$} \\ \hline
  & Soccer ball kicked  & 45     & 0.45 & \\ \hline
  & Car braking         &        & 1200 & 3.5 \\ \hline
  & Rocket launch       & 6\,000 &      & 30 \\ \hline
\end{tabular}
\end{ctab}

\taskpair{A net force of 24\,N acts on a 6\,kg mass. Find the acceleration.}{2.4cm}
         {A 0.5\,kg ball accelerates at 12\,m/s$^2$. Find the net force.}{2.4cm}

\taskpair{A force of 90\,N gives an object an acceleration of 3\,m/s$^2$. Find its mass.}{2.4cm}
         {A 1500\,kg car accelerates at 2.5\,m/s$^2$. Find the net force needed.}{2.4cm}

\taskpair{A 70\,kg sprinter accelerates at 4\,m/s$^2$ out of the blocks. Find the net force on her.}{2.4cm}
         {A 0.15\,kg cricket ball is hit with a net force of 30\,N. Find its acceleration.}{2.4cm}

\needspace{14\baselineskip}
\task A 2.5\,kg trolley is pushed with different net forces.
\pt{a}Complete the table.

\vspace{4pt}
\begin{ctab}
\begin{tabular}{|SZ{3.4cm}|Z{3.4cm}|Z{5.6cm}|}
  \hd & \thead{Net force (N)} & \thead{Mass (kg)} & \thead{Acceleration (m/s$^2$)} \\ \hline
  & 5  & 2.5 & \scaf{$a = 5 \div 2.5 = 2$} \\ \hline
  & 10 & 2.5 & \\ \hline
  & 15 & 2.5 & \\ \hline
  & 20 & 2.5 & \\ \hline
\end{tabular}
\end{ctab}
\pt{b}What happens to the acceleration each time the force goes up by 5\,N? Describe the pattern.
\answerbox{2.0cm}

\needspace{10\baselineskip}
\task A shopping trolley of mass 15\,kg is pushed with a force of 45\,N. Friction opposes the motion with 15\,N.
\pt{a}Find the net force.
\pt{b}Find the acceleration of the trolley.
\answerbox{3.0cm}

\needspace{10\baselineskip}
\task A 60\,kg cyclist and her 12\,kg bicycle accelerate at 1.5\,m/s$^2$.
\pt{a}Find the total mass.
\pt{b}Find the net force.
\pt{c}Air resistance is 30\,N. Find the force she must produce by pedalling.
\answerbox{3.6cm}

\needspace{10\baselineskip}
\task Two identical trolleys are pushed with the same force. Trolley A is empty, with a mass of 2\,kg. Trolley B carries an extra 6\,kg, so its total mass is 8\,kg.
\pt{a}Which trolley accelerates more?
\pt{b}Explain your answer using $F = ma$.
\pt{c}The force is 16\,N. Calculate both accelerations.
\answerbox{3.8cm}

\needspace{10\baselineskip}
\task A net force of 3600\,N gives a car an acceleration of 2\,m/s$^2$.
\pt{a}Find the mass of the car.
\pt{b}The car is then loaded with 200\,kg of luggage. Find its new acceleration with the same net force.
\answerbox{3.0cm}

\needspace{10\baselineskip}
\task A 900\,kg car travels along a flat road at a steady 60\,km/h. Its engine provides a driving force of 2000\,N.
\pt{a}What is the car's acceleration? Explain how you know.
\pt{b}What is the total size of the resistive forces (friction and air resistance) acting on the car?
\answerbox{3.0cm}

% ─── PART D ───────────────────────────────────────────────────────────────────
\lessonpart{D}{Above level}{Violet}\label{ext:5}

\begin{yr8box}
These questions go beyond the Year 7 standard. They ask you to combine ideas, use more than one formula, or judge someone else's reasoning.

\medskip
\needspace{10\baselineskip}
\exttask Ahmed pushes a 20\,kg crate with a force of 60\,N, but it does not move. He says: \emph{``$F = ma$ means it should accelerate at 3\,m/s$^2$, so the formula must be wrong.''} Explain Ahmed's mistake, and find the size of the friction force on the crate.
\answerbox{3.0cm}

\needspace{10\baselineskip}
\exttask Mia pushes a 50\,kg crate across the floor. Friction on the crate is 40\,N. She says: \emph{``If I push twice as hard, the crate will speed up twice as fast.''}
\pt{a}Find the acceleration when she pushes with 100\,N.
\pt{b}Find the acceleration when she pushes with 200\,N.
\pt{c}Is Mia right? Explain, using the idea of net force.
\answerbox{3.2cm}

\needspace{10\baselineskip}
\exttask A rocket's engines push exhaust gas downwards with a force of 500\,000\,N, and the gas pushes the rocket upwards with the same force. The rocket has a mass of 20\,000\,kg. Use $g = 9.8$\,N/kg.
\pt{a}Find the rocket's weight.
\pt{b}Find the net force on the rocket at lift-off.
\pt{c}Find its acceleration.
\pt{d}Which of Newton's laws explains (i) the upward push of the gas, and (ii) the size of the acceleration?
\answerbox{3.6cm}

\needspace{10\baselineskip}
\exttask A lift and its passengers have a total mass of 800\,kg. The cable pulls the lift upwards with a force of 8640\,N. Use $g = 9.8$\,N/kg.
\pt{a}Find the weight of the lift and its passengers.
\pt{b}Find the size and direction of the net force.
\pt{c}Find the acceleration of the lift.
\answerbox{2.8cm}

\needspace{10\baselineskip}
\exttask A cyclist speeds up from 2\,m/s to 8\,m/s in 4\,s. The cyclist and bicycle have a total mass of 80\,kg.
\pt{a}Calculate the acceleration, using $a = \Delta v \div t$ (change in speed $\div$ time).
\pt{b}Calculate the net force on the cyclist and bicycle.
\pt{c}Air resistance and friction total 40\,N. Find the force the cyclist must supply by pedalling.
\answerbox{2.8cm}

\needspace{10\baselineskip}
\exttask A 0.045\,kg golf ball is struck and leaves the tee at 60\,m/s. The club is in contact with the ball for 0.0005\,s.
\pt{a}Calculate the acceleration of the ball, using $a = \Delta v \div t$ (change in speed $\div$ time).
\pt{b}Calculate the force applied by the club.
\pt{c}Calculate the ball's weight, and comment on how the force from the club compares with it.
\answerbox{3.2cm}

\needspace{10\baselineskip}
\exttask In the trebuchet investigation at the end of this unit, a sling swings a 0.05\,kg projectile. While it is being swung, the net force on the projectile is 2\,N.
\pt{a}Find the projectile's acceleration.
\pt{b}The sling pulls on the projectile. Describe the reaction force.
\pt{c}Once the projectile leaves the sling, nothing pushes it forwards. Explain why it keeps moving forwards, and name the force that brings it back to the ground.
\answerbox{2.6cm}
\end{yr8box}

\newpage
% ══════════════════════════════════════════════════════════════════════════════
%  LESSON 6 — EXPERIMENT 2: INVESTIGATING NEWTON'S SECOND LAW      (Week 2 · P3)
% ══════════════════════════════════════════════════════════════════════════════
\sectiondiv{Experiment 2: Newton's Second Law}{Prediction \textbullet{} Variables \textbullet{} Trolley and pulley \textbullet{} Patterns in results}\label{sec:6}\label{exp:2}
\lessoncodes{VC2S8U14 \textbullet{} VC2S10U17 \textbullet{} VC2S8I02 \textbullet{} VC2S8I03 \textbullet{} VC2S8I04 \textbullet{} VC2S8I05}

% ─── PART A ───────────────────────────────────────────────────────────────────
\lessonpart{A}{Reading and comprehension}{Teal}\label{read:6}

\begin{readingbox}{Reading --- Slowing Down Gravity}
Drop a ball and it hits the floor in well under a second --- far too fast to measure by eye. Around 400 years ago, the Italian scientist Galileo Galilei faced exactly this problem. He wanted to find out how falling objects speed up, but he had no stopwatch, no camera and no clock accurate enough to help.

His clever solution was to slow everything down. Instead of dropping a ball, he let it roll down a long, gently sloping wooden ramp with a smooth groove cut into it. On a gentle slope, only a small part of the pull of gravity drives the ball along, so it speeds up slowly enough to measure. To time each roll, Galileo let water flow out of a large container through a thin pipe while the ball rolled, then weighed the water he had collected. More water meant more time.

Galileo's measurements showed that the ball sped up by the same amount every second --- it had a steady acceleration. Isaac Newton, born less than a year after Galileo died, built on this work to explain how force, mass and acceleration are linked.

In this experiment you will use the same trick. A trolley rolls along a flat runway, pulled by a small hanging mass on a string. Only the weight of that small mass pulls the trolley along, so it speeds up slowly enough for you to measure. By changing the hanging mass, you change the pulling force --- and you can test Newton's Second Law for yourself.
\end{readingbox}

\rcreset
\rcq{Why couldn't Galileo simply time a ball falling straight down?}{1.4cm}
\rcq{Explain how a gentle ramp helped Galileo, and how he measured time without a clock.}{2.0cm}
\rcq{In this experiment, what provides the force that pulls the trolley along?}{1.4cm}
\rcq{Why does using a \emph{small} hanging mass make the trolley's motion easier to measure?}{1.6cm}

% ─── PART B ───────────────────────────────────────────────────────────────────
\lessonpart{B}{Planning the experiment}{Blue}\label{theory:6}

\begin{expheader}{2}{Investigating Newton's Second Law}{To investigate how the size of the pulling force affects the acceleration of a trolley, when the total mass stays the same.}
\medskip\small\textbf{Curriculum:} VC2S8U14 \textbullet{} VC2S10U17 \textbullet{} VC2S8I02 \textbullet{} VC2S8I03 \textbullet{} VC2S8I04 \textbullet{} VC2S8I05
\end{expheader}

\subsection{The Set-up}\label{sub:6.1}

A trolley on a flat runway is joined by a string, over a pulley, to a hanging mass. The weight of the hanging mass provides the pulling force. The pulley simply turns the direction of the pull from \emph{down} to \emph{sideways}.

\diagExpTwoSetup
\expfigcap{Trolley connected by string over a pulley to a hanging mass. The hanging mass provides the driving force; the pulley redirects the tension.}

\needspace{24\baselineskip}
\subsection{The Forces on the Trolley}\label{sub:6.2}

Vertically, the trolley's weight and the normal force from the runway are \textbf{balanced}. Horizontally, the tension in the string is the only force that matters (friction is very small), so it is \textbf{unbalanced} --- and the trolley accelerates.

\diagExpTwoFBD
\expfigcap{Left: free body diagram of the trolley (tension drives it right; normal force and weight balance vertically). Right: free body diagram of the hanging mass (weight is larger than tension, so it accelerates downwards).}

\subsection{Calculating the Pulling Force}\label{sub:6.3}

The pulling force is the weight of the hanging mass, so use $W = mg$ with $g = 9.8$\,N/kg. The mass must be in \textbf{kilograms}: divide grams by 1000.

\begin{workedex}{Finding the pulling force}
A 150\,g mass hangs on the end of the string. Find the pulling force on the trolley.

\medskip
\textbf{Step 1 --- convert to kg:}\quad $150 \div 1000 = 0.150$\,kg

\textbf{Step 2 --- weight:}\quad $W = mg = 0.150 \times 9.8 = \mathbf{1.47\ \text{N}}$
\end{workedex}

\subsection{What $F = ma$ Predicts}\label{sub:6.4}

If $F = ma$, then doubling the force on the same mass should double the acceleration. This experiment tests that prediction.

\begin{scaffold}{Prediction --- before you begin}
What do you think will happen to the trolley's acceleration as the hanging mass is increased from 50\,g to 200\,g? Use $F = ma$ to give a reason.
\answerbox{2cm}
\end{scaffold}

\diagExpTwoForces
\expfigcap{Check your prediction against what Newton's Second Law predicts: as the hanging mass increases, the tension force increases, and so does the acceleration.}

\subsubsection*{Equipment}
Dynamics trolley \quad Runway \quad Pulley (clamped to the end of the bench) and string \quad Mass hanger and slotted masses (4 $\times$ 50\,g) \quad Motion sensor or ticker-tape timer \quad Balance \quad Buffer block and catch box

\subsubsection*{Method}
\begin{enumerate}[itemsep=2pt]
  \item Measure the mass of the trolley with the balance and record it.
  \item Clamp the pulley to the end of the bench. Tie a string from the trolley, over the pulley, to the mass hanger.
  \item Load the hanger so the hanging mass is \textbf{50\,g}. Put the other three 50\,g masses \textbf{on the trolley}.
  \item Hold the trolley still at the start line, then let go --- do not push it. Measure its acceleration with the motion sensor or ticker-tape timer.
  \item Repeat 3 times and record every result.
  \item Move one 50\,g mass from the trolley to the hanger (the hanging mass is now 100\,g) and repeat steps 4--5. Do the same for 150\,g and 200\,g.
\end{enumerate}

\begin{infobox}{Why move the masses instead of adding new ones?}
Moving a mass from the trolley to the hanger keeps the \textbf{total mass} being pulled the same. Only the \textbf{force} changes --- so it is a fair test.
\end{infobox}

\begin{infobox}{Safety}
Put a buffer block at the end of the runway and catch the trolley before it hits the pulley. Place a catch box under the hanging mass and keep your feet clear of it. Never overload the hanger.
\end{infobox}

% ─── PART C ───────────────────────────────────────────────────────────────────
\lessonpart{C}{Results and analysis}{Amber}\label{prac:6}

\subsubsection*{Results}

\begin{ctab}
\begin{tabular}{|SZ{2.2cm}|Z{2.6cm}|Z{2.1cm}|Z{2.1cm}|Z{2.1cm}|Z{2.5cm}|}
  \hd & \thead{Hanging mass (g)} & \thead{Pulling force $F = mg$ (N)} & \thead{Trial 1 (m/s$^2$)} & \thead{Trial 2 (m/s$^2$)} & \thead{Trial 3 (m/s$^2$)} & \thead{Mean accel. (m/s$^2$)} \\ \hline
  & \scaf{e.g.\ 50} & \scaf{0.49} & \scaf{0.38} & \scaf{0.42} & \scaf{0.41} & \scaf{0.40} \\ \hline
  & 50  & & & & & \\ \hline
  & 100 & & & & & \\ \hline
  & 150 & & & & & \\ \hline
  & 200 & & & & & \\ \hline
\end{tabular}
\end{ctab}

\begin{infobox}{Looking for a pattern}
Two quantities are \textbf{directly proportional} if doubling one doubles the other. To check, compare two rows: when the force doubles (50\,g $\rightarrow$ 100\,g), does the mean acceleration roughly double too?
\end{infobox}

\needspace{10\baselineskip}
\task Identify the independent variable, the dependent variable, and two controlled variables in this experiment.
\answerbox{2.6cm}

\needspace{10\baselineskip}
\task Why were the 50\,g masses moved from the trolley to the hanger, instead of adding new masses to the hanger?
\answerbox{2.2cm}

\needspace{10\baselineskip}
\task Describe the relationship between force and acceleration that you observed. Did doubling the force (50\,g $\rightarrow$ 100\,g) roughly double the acceleration? Was your prediction correct?
\answerbox{2.6cm}

\taskpair{A different hanging mass of 120\,g is used. Find the pulling force on the trolley.}{2.4cm}
         {Another group measured 0.81, 0.79 and 0.95\,m/s$^2$ for 100\,g. Which value looks wrong? Find the mean of the other two.}{2.4cm}

\needspace{10\baselineskip}
\task Use Newton's Second Law to explain \emph{why} a larger hanging mass gave the trolley a larger acceleration. Use the words \emph{net force} and \emph{mass}.
\answerbox{2.2cm}

\needspace{20\baselineskip}
\task Draw a line graph of your results. Put \textbf{pulling force (N)} on the horizontal axis and \textbf{mean acceleration (m/s$^2$)} on the vertical axis. Plot each point, draw a line of best fit, and give your graph a title.

\vspace{3pt}
\begin{tcolorbox}[enhanced, colback=white, colframe=MidGrey, boxrule=0.5pt, height=7.5cm,
  underlay={\begin{tcbclipinterior}\draw[step=5mm, Line, very thin] (interior.south west) grid (interior.north east);\end{tcbclipinterior}}]
\end{tcolorbox}

% ─── PART D ───────────────────────────────────────────────────────────────────
\lessonpart{D}{Above level}{Violet}\label{ext:6}

\begin{yr8box}
These questions go beyond the Year 7 standard. They ask you to interpret your graph, use a formula, or evaluate and extend the method.

\medskip
\needspace{10\baselineskip}
\exttask Look at your graph from Task~\ref{task:6.7}.
\pt{a}Is your line of best fit straight, and does it pass through (or close to) the origin? What does this tell you about force and acceleration?
\pt{b}Calculate the gradient of your line. What does it represent?
\answerbox{3.4cm}

\needspace{10\baselineskip}
\exttask For each row of your results, calculate $F \div a$ (pulling force $\div$ mean acceleration).
\pt{a}What do you notice about your answers?
\pt{b}Use $F = ma$ to explain what $F \div a$ should be equal to.
\pt{c}Compare your values with the total mass of the trolley plus all four 50\,g masses.
\answerbox{3.4cm}

\needspace{20\baselineskip}
\exttask Combining the two free body diagrams (Figure~\ref{fig:E6.2}) gives a formula for the acceleration of the whole system:

\diagExpTwoSystem
\expfigcap{Combining both free body diagrams gives the full system equation. When $M \gg m$, the result simplifies to $F = Ma$.}

In this experiment the total mass $M + m$ (the trolley plus all four slotted masses, wherever they sit) is 1.2\,kg.
\pt{a}Calculate the predicted acceleration when the hanging mass is 100\,g.
\pt{b}Your measured acceleration was probably a little smaller than this. Suggest one reason, and one change to the method that would reduce the difference.
\answerbox{3.4cm}

\needspace{12\baselineskip}
\exttask Design a second fair test, to find out how the \emph{mass} of the trolley affects its acceleration.
\pt{a}Name the independent and dependent variables.
\pt{b}What must you keep the same? Explain how you would keep the pulling force the same.
\pt{c}Predict what you would find, and give a reason using $F = ma$.
\answerbox{4.2cm}
\end{yr8box}

\end{document}
